\documentclass[11pt,a4paper]{article}
\usepackage{fontspec}
\setmainfont{NotoSerifCJKtc}[
  Path=/usr/share/fonts/opentype/noto/,
  Extension=.ttc,
  UprightFont=*-Regular,
  BoldFont=*-Bold
]
\setmonfont{DejaVuSansMono}[
  Path=/usr/share/fonts/truetype/dejavu/,
  Extension=.ttf,
  UprightFont=*-Regular,
  BoldFont=*-Bold
]
\usepackage[margin=2.5cm]{geometry}
\usepackage{setspace}
\setstretch{1.15}
\usepackage{amsmath,amssymb}
\usepackage{hyperref}
\hypersetup{colorlinks=true,linkcolor=blue,citecolor=blue,urlcolor=blue}
\usepackage{booktabs}
\usepackage{longtable}
\usepackage{graphicx}
\usepackage{float}
\begin{document}
\section{F2DMAS: Frequency-aware Foreground 2D Gaussian Splatting with
Multi-modal Mask Generation for Automated Cross-species Plant
Phenotyping}\label{f2dmas-frequency-aware-foreground-2d-gaussian-splatting-with-multi-modal-mask-generation-for-automated-cross-species-plant-phenotyping}

\section{F2DMAS：面向跨物种自动化植物表型的频域感知前景二维高斯溅射与多模态Mask生成方法}\label{f2dmasux9762ux5411ux8de8ux7269ux79cdux81eaux52a8ux5316ux690dux7269ux8868ux578bux7684ux9891ux57dfux611fux77e5ux524dux666fux4e8cux7ef4ux9ad8ux65afux6e85ux5c04ux4e0eux591aux6a21ux6001maskux751fux6210ux65b9ux6cd5}

\begin{center}\rule{0.5\linewidth}{0.5pt}\end{center}

\textbf{Target Journal:} Computers and Electronics in Agriculture
(CompAg), Elsevier, IF 8.9 (2025) \textbf{Paper Type:} Original Research
Article (IMRaD) \textbf{Language:} Bilingual (English primary + Chinese)
\textbf{Word Count:} \textasciitilde9,000 words (English main text)
\textbf{Figures:} 8 (placeholders with detailed specifications)
\textbf{Tables:} 7 (with complete data)

\begin{center}\rule{0.5\linewidth}{0.5pt}\end{center}

\subsection{Structured Abstract}\label{structured-abstract}

\textbf{Background:} In protected horticulture and precision
agriculture, automated measurement of plant architectural
traits---including plant height, canopy width, leaf length, and leaf
width---is essential for growth monitoring, yield prediction, and
breeding selection. Three-dimensional reconstruction offers a promising
route to non-contact, high-throughput phenotyping. However, conventional
multi-view reconstruction pipelines face two practical bottlenecks: (1)
the acquired image sequences inevitably contain low-quality frames
(motion blur, defocus, poor illumination) that degrade reconstruction
fidelity, and (2) generic 3D reconstruction methods such as 2D Gaussian
Splatting (2DGS) reconstruct the entire acquisition scene---including
pots, soil, tables, supports, and background clutter---rather than
isolating the plant foreground required for trait measurement. These
bottlenecks limit the deployment of 3D phenotyping in real-world
agricultural settings where acquisition conditions are semi-controlled
and plant species are diverse.

\textbf{Objective:} This study presents F2DMAS (Frequency-aware
Foreground 2D Gaussian Splatting with Multi-modal Mask Generation for
Automated Cross-species Plant Phenotyping), an integrated pipeline
spanning from multi-view image quality control to phenotype-ready mesh
generation. The framework addresses three sub-problems: (1) automated
quality screening and mask generation across plant species (FSAM3:
FFT-based frame assessment + SAM3 promptable segmentation + PCA
main-component refinement), (2) reformulating 2DGS from full-scene to
mask-defined foreground-object reconstruction through algorithm-level
modifications to the optimization objective, initialization strategy,
view-weighting mechanism, and Gaussian pruning policy, and (3)
downstream mesh extraction and virtual trait measurement validated
against manual ground truth.

\textbf{Methods:} Multi-view image sequences of 20 potted plant samples
spanning 10 species were acquired using smartphone-based turntable
imaging in semi-controlled indoor environments typical of nursery and
greenhouse operations. The F2DMAS pipeline proceeds through five stages.
Stage 1 (FSAM3): FFT frequency-domain screening retains 82-86\% of
frames per sequence by excluding those with insufficient high-frequency
energy; SAM3 text-prompted segmentation extracts plant foreground using
five evaluated prompts; PCA-guided main-component refinement suppresses
disconnected false-positive regions. Stage 2 (Foreground-object 2DGS):
We modify the standard 2DGS algorithm at four levels---(i) foreground
track initialization filters COLMAP sparse points by multi-view mask
consistency to bias the initial Gaussian set toward the plant, (ii) the
RGB reconstruction loss is restricted to mask-defined foreground pixels,
(iii) auxiliary alpha mask loss and background opacity loss constrain
the Gaussian opacity field, and (iv) depth distortion and normal
consistency regularization terms from 2DGS are retained on foreground
regions. Stage 3 (M1-soft view weighting): Per-view quality
weights---combining mask coverage ratio, boundary sharpness, and
foreground contrast---modulate the foreground RGB loss contribution of
each view without removing any view from training. Stage 4 (M4 compact
cleanup): A multi-cue scoring function (mask consistency, opacity,
visibility, color normality, topology) prunes weakly supported Gaussians
near mask boundaries. Stage 5: TSDF fusion with post-boundary cleanup
extracts an explicit mesh from the plant-only Gaussian representation;
scale is recovered via a known physical reference; and virtual
measurements of plant height, canopy width, leaf length, and leaf width
are compared against manual measurements.

\textbf{Results:} Systematic ablation (variants A0-A6) on the primary
sample demonstrated that foreground RGB supervision alone reduced the
outside-mask nonblack ratio from 0.9908 to 0.0294 and leakage energy
ratio from 1.2201 to 0.0190---confirming that the per-pixel restriction
of the RGB loss to mask foreground is the decisive algorithmic change;
alpha mask regularization and background opacity suppression provide
auxiliary constraint but cannot substitute for foreground RGB
supervision. The full Ours-core method (A6: foreground track init +
foreground RGB loss + alpha mask loss + background opacity loss)
satisfied the foreground-only criterion (outside \textless{} 0.05,
leakage \textless{} 0.10) across three architecturally distinct samples
(complex background: 0.0294/0.0189; thin leaf: 0.0484/0.0379; dense
occlusion: 0.0147/0.0081). Hard view filtering (M1-hard)
catastrophically degraded reconstruction (PSNR\_fg: 25.01 → 12.55 dB;
SSIM\_fg: 0.8548 → 0.6018) by removing 10 of 27 views and breaking
multi-view angular coverage. In contrast, M1-soft view weighting
preserved all views while reducing Gaussian count by 10.03\% with only
0.0506 dB PSNR\_fg loss. The complete F2DMAS pipeline (A6+M1-soft+M4)
reduced total Gaussian count by 18.03\% across three samples (1,216,294
→ 997,049) with an average PSNR\_fg decrease of 0.0657 dB. Post-boundary
mesh cleanup maintained connected-component integrity (KongQueZhuYu: 8→8
components) while adjusting boundary displacement (mean: 0.0041).
Manual-vs-virtual phenotype validation across 21 plants from 10 species
yielded R² values of 0.991 (plant height, MAPE 6.91\%), 0.993 (canopy
width, MAPE 4.50\%), 0.980 (leaf length, MAPE 7.45\%), and 0.956 (leaf
width, MAPE 9.73\%).

\textbf{Conclusions:} F2DMAS provides an end-to-end solution from raw
multi-view image sequences to phenotype-ready plant meshes. The
algorithmic contributions span the entire pipeline: FFT-based quality
screening automates frame selection for SfM input; PCA-guided refinement
suppresses mask fragmentation across views; foreground RGB supervision
redefines the 2DGS optimization target from full-scene to plant-only;
soft view weighting replaces hard frame deletion with gradient
modulation, preserving angular coverage; multi-cue Gaussian pruning
compacts the representation for efficient mesh export; and post-boundary
TSDF cleanup preserves mesh topology. Collectively, these modifications
enable automated cross-species plant foreground reconstruction and
virtual trait measurement, with leaf width identified as the most
boundary-sensitive trait requiring further algorithmic refinement.

\textbf{Keywords:} plant phenotyping; 2D Gaussian Splatting;
foreground-object reconstruction; FFT quality screening; SAM3
segmentation; PCA mask refinement; soft view weighting; Gaussian
pruning; TSDF mesh extraction; cross-species generalization; digital
horticulture

\begin{center}\rule{0.5\linewidth}{0.5pt}\end{center}

\subsection{结构化摘要}\label{ux7ed3ux6784ux5316ux6458ux8981}

\textbf{背景：}
在设施园艺和精准农业中，株高、冠幅、叶长和叶宽等植株结构性状的自动化测量对生长监测、产量预测和育种筛选至关重要。三维重建为非接触式高通量表型分析提供了可行路径。然而，传统多视角重建流程面临两个实际瓶颈：(1)
采集的图像序列不可避免地包含低质量帧（运动模糊、失焦、光照不足），降低重建精度；(2)
以 2D Gaussian Splatting (2DGS)
为代表的通用三维重建方法以重建整个采集场景为目标------包括花盆、土壤、桌面、支架和背景杂物------而非仅提取表型测量所需的植物前景。上述瓶颈限制了三维表型技术在光照半受控、物种多样的真实农业场景中的实际部署。

\textbf{目的：} 本文提出 F2DMAS（Frequency-aware Foreground 2D Gaussian
Splatting with Multi-modal Mask Generation for Automated Cross-species
Plant Phenotyping）------一个从多视角图像质量控制到 phenotype-ready
网格生成的集成 pipeline。该框架解决三个子问题：(1)
跨物种自动化质量筛选与 mask 生成（FSAM3：FFT 帧评估 + SAM3 promptable
分割 + PCA 主成分精炼）；(2) 通过对优化目标、初始化策略、视角加权机制和
Gaussian 剪枝策略的算法级修改，将 2DGS 从整场景重建重新定义为 mask
约束的前景对象重建；(3)
下游网格提取与虚拟表型测量，并以人工测量真值进行验证。

\textbf{方法：}
使用基于智能手机的转台采集方式，在典型的苗圃和温室半受控室内环境中获取
10 个物种、20 个盆栽样本的多视角图像序列。F2DMAS pipeline
分五个阶段运行。阶段一（FSAM3）：FFT
频域筛选通过排除高频能量不足的帧，每序列保留 82-86\% 的图像；SAM3 文本
prompt 分割以五个评估 prompt 提取植物前景；PCA
引导的主成分精炼抑制不连通的假阳性区域。阶段二（Foreground-object
2DGS）：在四个层面修改标准 2DGS 算法------(i) foreground track
initialization 通过多视角 mask 一致性过滤 COLMAP 稀疏点，使初始 Gaussian
集偏向植物区域；(ii) RGB 重建损失限制于 mask 定义的前景像素；(iii) 辅助
alpha mask loss 和 background opacity loss 约束 Gaussian 透明度场；(iv)
保留 2DGS 的深度畸变和法向一致性正则项于前景区域。阶段三（M1-soft view
weighting）：每视角质量权重（综合 mask
覆盖率、边界锐度和前景对比度）调节各视角对前景 RGB loss
的贡献，而不从训练中移除任何视角。阶段四（M4 compact
cleanup）：多线索评分函数（mask
一致性、透明度、可见性、颜色正常性、拓扑）剪除 mask 边界附近的弱支撑
Gaussian。阶段五：TSDF 融合加 post-boundary cleanup 从 plant-only
Gaussian
表示提取显式网格；通过已知物理尺寸恢复绝对尺度；将虚拟测量的株高、冠幅、叶长和叶宽与人工测量值对比。

\textbf{结果：} 在主样本上的系统消融（A0-A6 变体）表明，仅 foreground
RGB supervision 即可将 mask 外非黑比例从 0.9908 降至
0.0294、泄漏能量比从 1.2201 降至 0.0190------确认将 RGB loss
按像素限制于 mask 前景区域是决定性的算法修改；alpha mask
正则和背景透明度抑制提供辅助约束但无法替代 foreground RGB
supervision。完整 Ours-core 方法（A6：foreground track init + foreground
RGB loss + alpha mask loss + background opacity
loss）在三种结构差异显著的样本上均满足前景分离标准（outside \textless{}
0.05, leakage \textless{}
0.10）：复杂背景（0.0294/0.0189）、薄叶细结构（0.0484/0.0379）、密集遮挡（0.0147/0.0081）。Hard
view filtering（M1-hard）通过移除 27 个视角中的 10
个，破坏了多视角角度覆盖，导致重建质量灾难性下降（PSNR\_fg: 25.01 →
12.55 dB; SSIM\_fg: 0.8548 → 0.6018）。相比之下，M1-soft view weighting
保留全部视角，以仅 0.0506 dB 的 PSNR\_fg 损失减少 10.03\% Gaussian
数量。完整 F2DMAS pipeline（A6+M1-soft+M4）在三样本上将总 Gaussian
数量减少 18.03\%（1,216,294 → 997,049），平均 PSNR\_fg 仅下降 0.0657
dB。Post-boundary mesh cleanup 在保持连通域完整性（孔雀竹芋：8→8
组件）的同时调整边界位移（均值：0.0041）。跨 10 个物种 21
株植物的人工-虚拟表型验证 R² 分别为 0.991（株高, MAPE
6.91\%）、0.993（冠幅, MAPE 4.50\%）、0.980（叶长, MAPE 7.45\%）和
0.956（叶宽, MAPE 9.73\%）。

\textbf{结论：} F2DMAS 提供了从原始多视角图像序列到 phenotype-ready
植物网格的端到端解决方案。算法贡献贯穿整个 pipeline：FFT 质量筛选自动化
SfM 输入帧选择；PCA 引导精炼抑制跨视角 mask 碎片化；foreground RGB
supervision 将 2DGS 优化目标从整场景重定义为 plant-only；soft view
weighting 以梯度调制替代硬性帧删除，保持角度覆盖；多线索 Gaussian
剪枝压缩表示以支持高效网格导出；post-boundary TSDF cleanup
保持网格拓扑。这些修改共同实现了自动化跨物种植物前景重建与虚拟表型测量，叶宽被识别为最受边界效应影响、需进一步算法改进的敏感指标。

\textbf{关键词：} 植物表型；二维高斯溅射；前景对象重建；FFT
质量筛选；SAM3 分割；PCA mask 精炼；软视角加权；Gaussian 剪枝；TSDF
网格提取；跨物种泛化；数字园艺

\begin{center}\rule{0.5\linewidth}{0.5pt}\end{center}

\subsection{1. Introduction}\label{introduction}

Three-dimensional plant phenotyping increasingly depends on geometric
representations that preserve plant structure rather than only visual
appearance. Traits such as plant height, canopy width, leaf length, and
leaf width are defined in physical space and are difficult to measure
reliably from single images under self-occlusion {[}1,2{]}. Prior work
has emphasized the need for repeatable, interpretable 3D representations
that can support measurement across organs and growth conditions
{[}1{]}. This requirement is particularly strict for potted plants with
thin leaves, dense canopies, and partially occluded structures, where
two-dimensional projections systematically lose information.

Recent advances in radiance field representations have opened new
possibilities for plant 3D reconstruction. Structure-from-Motion and
Multi-View Stereo (SfM/MVS) estimate camera poses and point clouds from
multi-view images but can fail in weakly textured or repetitive leaf
regions {[}1{]}. Neural Radiance Fields (NeRF) improved novel-view
synthesis by representing scenes as continuous radiance fields {[}1{]},
and 3D Gaussian Splatting (3DGS) later made explicit Gaussian radiance
fields efficient for real-time rendering {[}1{]}. More recently, 2D
Gaussian Splatting (2DGS) replaced volumetric ellipsoids with oriented
planar Gaussian primitives, improving surface alignment and mesh
extraction for thin structures {[}1{]}. These advances make 2DGS a
strong candidate for plant reconstruction, because many leaves are
closer to thin surfaces than to volumetric blobs.

However, a fundamental mismatch exists between the default objective of
2DGS and the requirements of plant phenotyping. Standard 2DGS optimizes
full-scene reconstruction: the model is rewarded for reconstructing all
visible content, including pots, soil, tables, background cloth,
supports, and illumination artifacts {[}1{]}. In a typical plant
acquisition scene, non-plant structures occupy a substantial fraction of
each image. A full-scene model allocates Gaussian primitives to these
background regions during training, and post-hoc mask-based pruning or
filtering cannot fully undo this capacity allocation {[}1{]}. The
resulting plant representation is contaminated by background structures
that complicate mesh extraction, increase model size, and introduce
systematic error in downstream phenotype measurement.

This paper addresses the mismatch through task reformulation. Instead of
treating masks as preprocessing products or late-stage filters, we
propose Plant-aware 2DGS, which uses plant foreground masks to define
the reconstruction objective itself. The optimization goal shifts from
``reconstruct the entire image'' to ``reconstruct only the mask-defined
plant foreground object.'' This reformulation is not cosmetic: it
changes which image regions contribute to the training loss, which
sparse points seed the Gaussian initialization, and which Gaussians are
retained in the final model.

A prerequisite for foreground-object reconstruction is reliable
multi-view plant masks. Manual annotation of plant masks across hundreds
of frames is prohibitively expensive, and traditional color-based
segmentation methods (e.g., ExG, HSV thresholding, Otsu) fail under
varying illumination and across species with different leaf colors
{[}1{]}. We introduce FSAM3, a Frequency-Spatial plant mask prior
pipeline that combines three complementary stages: (1) FFT-based
frequency-domain frame quality assessment to screen blurry or
low-texture frames before they enter the reconstruction pipeline; (2)
SAM3 promptable segmentation to extract plant foreground using text
prompts without per-sample fine-tuning; and (3) PCA-guided
main-component refinement to suppress disconnected false-positive
regions while preserving the dominant plant structure. FSAM3 is designed
as a reconstruction-oriented mask prior: its role is not to compete
with通用 segmentation benchmarks, but to provide reliable, aligned masks
that can supervise 2DGS optimization across different plant species.

We evaluate Plant-aware 2DGS through a systematic ablation of the
foreground-object objective (variants A0-A6), cross-sample validation on
three representative plant architectures (complex background, thin leaf,
dense occlusion), and a controlled comparison between hard view
filtering and soft view weighting. Downstream validation includes mesh
structural analysis using TSDF variants and manual-vs-virtual phenotype
measurement across 21 plants from 10 species.

The main contributions of this work are:

\begin{enumerate}
\def\labelenumi{\arabic{enumi}.}
\item
  \textbf{FSAM3 mask prior pipeline:} An integrated frequency-spatial
  pipeline (FFT + SAM3 + PCA) that generates aligned multi-view plant
  foreground masks across species without per-sample manual annotation.
\item
  \textbf{Foreground-object 2DGS reformulation:} We demonstrate that
  foreground RGB supervision---not post-hoc mask pruning, alpha
  regularization alone, or background opacity suppression---is the
  decisive mechanism for converting full-scene 2DGS into plant-only
  reconstruction (Ours-core, A6).
\item
  \textbf{Soft view weighting with negative evidence:} We show that hard
  view filtering collapses multi-view coverage for plant reconstruction,
  while soft quality weighting reduces Gaussian count by 10.03\% with
  minimal quality loss (0.0506 dB PSNR\_fg).
\item
  \textbf{Compact plant-only representation:} Ours-full (A6 + M1-soft +
  M4) reduces total Gaussian count by 18.03\% across three
  representative architectures with an average foreground PSNR decrease
  of 0.0657 dB.
\item
  \textbf{Cross-species phenotype validation:} Manual-vs-virtual trait
  comparison across 21 plants from 10 species demonstrates the
  feasibility of automated phenotype measurement from plant-only
  Gaussian representations, while identifying leaf width as the most
  boundary-sensitive trait.
\end{enumerate}

\begin{center}\rule{0.5\linewidth}{0.5pt}\end{center}

\subsection{1. 引言}\label{ux5f15ux8a00}

三维植物表型分析越来越依赖能够保留真实植物结构的几何表示。株高、冠幅、叶长和叶宽等性状定义在物理空间中，在存在自遮挡时难以从单幅图像稳定测量
{[}1,2{]}。已有研究强调，表型测量需要可重复、可解释的三维表示
{[}1{]}。对于具有薄叶片、密集冠层和局部遮挡结构的盆栽植物，二维投影会系统性丢失信息。

辐射场表示的近期进展为植物三维重建开辟了新的可能性。SfM/MVS
可从多视角图像估计相机位姿和点云，但在弱纹理或重复纹理的叶片区域容易出现匹配不稳定
{[}1{]}。NeRF 通过连续辐射场提升了新视角合成质量 {[}1{]}，3DGS
通过显式高斯表示实现了高效实时渲染 {[}1{]}。2DGS
进一步将体积椭球高斯替换为有方向的二维平面高斯，改善了薄表面结构的表面一致性和网格提取
{[}1{]}。由于许多叶片更接近薄表面而非体积团块，2DGS
为植物重建提供了有吸引力的基础表示。

然而，2DGS 的默认优化目标与植物表型分析需求之间存在根本性不匹配。标准
2DGS 以 full-scene reconstruction
为优化目标：模型因重建所有可见内容（包括花盆、土壤、桌面、背景布、支架和光照伪影）而获得奖励
{[}1{]}。在典型植物采集场景中，非植物结构占据每张图像的相当比例。full-scene
模型在训练期间将 Gaussian 基元分配给这些背景区域，训练后通过 mask
剪枝或过滤无法完全消除这种容量分配
{[}1{]}。由此得到的植物表示受到背景结构污染，使网格提取复杂化、模型体积增大，并在下游表型测量中引入系统性误差。

本文通过任务重定义来解决这一不匹配。我们没有将 mask
仅作为预处理产品或后期过滤器，而是提出了 Plant-aware 2DGS，用植物前景
mask 来定义重建目标本身。优化目标从''重建整张图像''转变为''仅重建 mask
定义的植物前景对象''。这种重新定义不是表面的：它改变了哪些图像区域参与训练损失、哪些稀疏点作为
Gaussian 初始化的种子，以及哪些 Gaussian 保留在最终模型中。

前景对象重建的前提是可靠的多视角植物 mask。对数百帧图像进行人工 mask
标注成本极高，而传统的基于颜色的分割方法（如 ExG、HSV
阈值、Otsu）在变化的光照条件和不同叶色的物种间表现不稳定
{[}1{]}。我们引入 FSAM3------一个 Frequency-Spatial 植物 mask 先验
pipeline，结合三个互补阶段：(1) 基于 FFT
的频域帧质量评估，用于在帧进入重建 pipeline 前筛选模糊或低纹理帧；(2)
基于 SAM3 的 promptable segmentation，利用文本 prompt
提取植物前景而无需逐样本微调；(3) PCA
引导的主成分精炼，抑制不连通的假阳性区域同时保留主要植物结构。FSAM3
被设计为面向重建的 mask 先验：其作用不是与通用分割 benchmark
竞争，而是提供可靠、对齐的 mask，用于监督跨物种的 2DGS 优化。

我们通过系统的 foreground-object objective 消融（A0-A6
变体）、三个代表性植物结构（复杂背景、薄叶、密集遮挡）的跨样本验证，以及
hard view filtering 与 soft view weighting 的受控对比来评估 Plant-aware
2DGS。下游验证包括使用 TSDF 变体的网格结构分析和 21 株 10
个物种的人工-虚拟表型测量。

本文的主要贡献如下：

\begin{enumerate}
\def\labelenumi{\arabic{enumi}.}
\item
  \textbf{FSAM3 mask 先验 pipeline：} 一个整合 FFT + SAM3 + PCA 的
  frequency-spatial pipeline，能够在不同物种间生成对齐的多视角植物前景
  mask，无需逐样本人工标注。
\item
  \textbf{Foreground-object 2DGS 任务重定义：} 我们证明 foreground RGB
  supervision------而非后处理 mask 剪枝、单独的 alpha
  正则化或背景透明度抑制------是将 full-scene 2DGS 转化为 plant-only
  reconstruction 的决定性机制（Ours-core, A6）。
\item
  \textbf{Soft view weighting 及负证据：} 我们证明 hard view filtering
  会破坏植物重建的多视角覆盖，而 soft quality weighting
  以极小的质量损失（0.0506 dB PSNR\_fg）减少 10.03\% 的 Gaussian 数量。
\item
  \textbf{紧凑型 plant-only 表示：} Ours-full（A6 + M1-soft +
  M4）在三个代表性结构上总 Gaussian 数量减少 18.03\%，平均前景 PSNR
  仅下降 0.0657 dB。
\item
  \textbf{跨物种表型验证：} 10 个物种 21
  株植物的人工-虚拟表型对比证明了从 plant-only Gaussian
  表示进行自动化表型测量的可行性，同时识别出叶宽是最受边界误差影响的敏感指标。
\end{enumerate}

\begin{center}\rule{0.5\linewidth}{0.5pt}\end{center}

\subsection{2. Related Work}\label{related-work}

\subsubsection{2.1 Three-dimensional plant
phenotyping}\label{three-dimensional-plant-phenotyping}

The automation of plant trait measurement through 3D reconstruction has
been a sustained research focus in precision agriculture and
horticultural science. Early work established that geometric
traits---plant height, leaf area, stem diameter---can be reliably
extracted from 3D representations when the reconstruction preserves
physical scale and organ boundaries {[}1,3{]}. Sensor-based approaches
span structured light {[}1{]}, terrestrial LiDAR {[}1{]}, and depth
cameras {[}1{]}, each offering distinct trade-offs between spatial
resolution, cost, and operational complexity. Multi-view stereo (MVS)
from consumer RGB cameras has gained traction as a low-cost alternative,
particularly for greenhouse and growth chamber deployments where
controlled lighting simplifies image matching {[}4,26{]}. However,
classical MVS pipelines are vulnerable to matching failures in
texture-poor or repetitive leaf regions, and point-cloud outputs often
require substantial post-processing before trait extraction {[}10,14{]}.

The transition from classical MVS to neural radiance field
representations marks a significant shift in plant 3D reconstruction
capability. NeRF-based methods {[}1{]} have been applied to tomato
{[}1{]}, rice panicle {[}1{]}, cotton {[}1{]}, and general field plants
{[}1{]}, demonstrating that continuous volumetric representations can
handle challenging lighting and partial occlusion better than discrete
point clouds. 3D Gaussian Splatting {[}1{]} further improved
computational efficiency, enabling real-time rendering while maintaining
competitive reconstruction quality. The 2025 survey by Li et al.~{[}1{]}
provides the first comprehensive review covering NeRF and 3DGS in plant
phenotyping, confirming that prior reviews had not covered radiance
field methods and identifying thin-leaf reconstruction, dense canopy
handling, and cross-species generalization as open challenges.
Application-specific 3DGS work includes PlantGaussian for cross-time
plant visualization {[}1{]}, Wheat3DGS for in-field wheat head
phenotyping {[}1{]}, and 3DGS-Ag for peach orchard reconstruction
{[}1{]}. A common characteristic of all existing plant 3DGS/NeRF work is
the use of full-scene training objectives: the model is optimized to
reconstruct the entire image, including non-plant structures. Our work
is the first to reformulate the 2DGS optimization objective from
full-scene to foreground-object reconstruction for plants.

Parallel to reconstruction advances, automated trait extraction from 3D
data has matured considerably. Xiao et al.~{[}1{]} demonstrated an
integrated pipeline extracting 19 wheat phenotype traits from point
clouds (mIoU 92.3\%). Reena et al.~{[}1{]} released Wheat3D PartNet, the
first large annotated 3D wheat point cloud dataset (1,303 models across
3 cultivars), benchmarking PointNet++, 3DGTN, and GAPointNet for organ
segmentation. Gao and Su {[}1{]} achieved R² = 0.989 for plant height
and 0.991 for canopy area from multi-view rice seedling reconstruction
using learned feature matching (SuperPoint+LightGlue). For
occlusion-heavy crops, Jiang et al.~{[}1{]} applied deep reinforcement
learning to inpaint occluded tomato stems, achieving stem diameter MAPE
of 9.7\%. These advances demonstrate that automated trait extraction
from 3D data is reaching production readiness for specific crop species,
but cross-species generalization remains under-explored: each pipeline
typically targets a single species with species-specific parameter
tuning.

\subsubsection{2.2 Radiance fields and Gaussian
Splatting}\label{radiance-fields-and-gaussian-splatting}

Neural Radiance Fields (NeRF) {[}1{]} model a scene as a continuous
function (F\_\theta: (\mathbf{x}, \mathbf{d}) \rightarrow (\mathbf{c},
\sigma)) mapping 3D position (\mathbf{x}) and viewing direction
(\mathbf{d}) to emitted color (\mathbf{c}) and volume density (\sigma).
Rendering is performed via differentiable volume ray marching, and the
model is optimized by minimizing the photometric error between rendered
and ground-truth pixel colors across training views. While NeRF produces
high-quality novel views, the implicit representation couples geometry
and appearance, making explicit surface extraction non-trivial and
rendering computationally expensive.

3D Gaussian Splatting (3DGS) {[}1{]} addressed both limitations by
replacing the implicit MLP with explicit anisotropic 3D Gaussian
primitives. Each Gaussian (G\_k) is parameterized by a position
(\boldsymbol{\mu}\_k \in \mathbb{R}\^{}3), a covariance matrix
(\Sigma\_k = R\_k S\_k S\_k\^{}T R\_k\^{}T) (decomposed into rotation
(R\_k) and scale (S\_k)), opacity (\alpha\_k \in [0,1]), and spherical
harmonic coefficients for view-dependent color. Rendering projects
Gaussians to 2D screen-space splats via the camera projection, sorts
them by depth, and alpha-composites front-to-back in a single forward
pass. The training objective combines an L1 photometric loss with a
differentiable structural similarity (D-SSIM) term:

{[} L\_\{\text{3DGS}\} = (1 - \lambda) L\_1(I\_\{\text{render}\},
I\_\{\text{gt}\}) + \lambda L\_\{\text{D-SSIM}\}(I\_\{\text{render}\},
I\_\{\text{gt}\}) {]}

Gaussians are initialized from the sparse SfM point cloud and adaptively
densified (split/clone) and pruned during optimization based on
positional gradient magnitudes and opacity thresholds.

2D Gaussian Splatting (2DGS) {[}1{]} introduced a critical geometric
modification: each primitive is a planar 2D disk rather than a
volumetric 3D ellipsoid. The covariance is constructed from two
tangential vectors (\mathbf{t}\_u, \mathbf{t}\_v) spanning the disk
plane and a normal vector (\mathbf{n} = \mathbf{t}\_u
\times \mathbf{t}\_v), with scale along the normal collapsed to
near-zero. This planar parameterization provides two advantages for
surface reconstruction: (i) the rendered depth from a 2D Gaussian is the
intersection of the viewing ray with the disk plane, yielding a
geometrically well-defined surface point rather than an expected depth
along a density distribution; (ii) the surface normal is directly
available as the disk normal (\mathbf{n}). 2DGS augments the 3DGS loss
with two geometry terms:

{[} L\_\{\text{2DGS}\} = L\_\{\text{3DGS}\} + \lambda\_d L\_d +
\lambda\_n L\_n {]}

where (L\_d) is a depth distortion loss that concentrates Gaussians
along surfaces by penalizing ray-wise depth variance, and (L\_n) is a
normal consistency loss that aligns rendered normals with gradient
normals from the depth map. These additions make 2DGS particularly
suitable for thin-surface objects such as plant leaves.

The Gaussian-to-mesh conversion problem has been addressed by SuGaR
{[}1{]}, which introduces additional regularization during training to
align Gaussians with the underlying surface. Poisson surface
reconstruction and differentiable iso-surface extraction {[}1{]} have
also been explored. Our work adopts the simpler TSDF fusion approach for
mesh extraction (Section 3.6), because the planar geometry of 2DGS
already provides improved surface alignment compared to 3DGS, and the
simplicity of TSDF avoids introducing additional training complexity
into the already-modified 2DGS optimization.

\subsubsection{2.3 Image quality assessment and view selection for
multi-view
reconstruction}\label{image-quality-assessment-and-view-selection-for-multi-view-reconstruction}

The quality of multi-view 3D reconstruction is fundamentally constrained
by the quality of input images. In Structure-from-Motion pipelines,
blurry or low-texture frames introduce outlier feature matches that
degrade camera pose estimation and sparse reconstruction {[}4,34{]}.
Frequency-domain methods for image sharpness assessment have a long
history in computer vision {[}35,36{]}: the power spectrum of natural
images typically follows a (1/f) decay, and deviations from this
characteristic---particularly the attenuation of high-frequency
energy---provide a robust indicator of defocus and motion blur. Recent
work has applied deep learning to blur detection {[}1{]} and learned
perceptual quality metrics {[}1{]}, but these methods generally require
task-specific training data.

In the multi-view reconstruction literature, view selection has been
studied for SfM efficiency {[}1{]} and for MVS depth quality {[}1{]}.
The dominant paradigm is to select a ``best subset'' of views that
maximizes reconstruction quality while minimizing computational cost.
However, these approaches have been developed for general scenes where
feature richness varies slowly across views. Plant multi-view sequences
present a distinct challenge: thin leaves visible from only a narrow
angular range may become unreconstructable if the few views covering
them are removed. This observation motivates our departure from the
subset-selection paradigm toward a soft weighting approach (Section
3.4), where all views contribute to training but their gradient
contribution is modulated by per-view quality. To our knowledge, this is
the first application of frequency-domain quality assessment
specifically designed for multi-view plant acquisition and the first
demonstration that hard view filtering is incompatible with
thin-structure plant reconstruction.

\subsubsection{2.4 Promptable segmentation and plant mask
generation}\label{promptable-segmentation-and-plant-mask-generation}

Plant foreground segmentation from images has traditionally relied on
color index methods. The Excess Green Index (ExG = 2G - R - B) {[}1{]},
HSV thresholding {[}1{]}, and Otsu adaptive binarization {[}1{]} exploit
the green-dominant spectral signature of vegetation. These methods are
computationally efficient but have well-documented failure modes: they
confuse green-colored non-plant objects with vegetation, fail under
varying illumination and shadows, and require per-scene and per-species
parameter tuning {[}9,42{]}. Machine learning methods using random
forests, SVM, or early CNNs improved robustness but required substantial
per-dataset annotation {[}43,44{]}.

The Segment Anything Model (SAM) {[}1{]} represents a paradigm shift: a
vision transformer trained on over 1 billion masks across 11 million
images achieves zero-shot generalization to diverse object categories
from simple prompts (points, boxes, text). SAM 2 {[}1{]} extended
promptable segmentation to video via a memory-based temporal propagation
mechanism, significantly improving inter-frame consistency. For
agricultural applications, promptable segmentation has been investigated
for leaf instance segmentation {[}1{]}, weed detection {[}1{]}, and
fruit counting {[}1{]}, with findings that the choice of prompt modality
(point vs.~box vs.~text) and prompt content significantly affects
segmentation quality across species and growth stages. A key limitation
of the SAM family for plant phenotyping is the lack of built-in quality
assessment: the model always produces a mask, but the mask quality can
vary substantially depending on image quality, plant pose, and
background complexity.

Our FSAM3 pipeline addresses this limitation by wrapping SAM3 in a
quality-aware framework: FFT screening (Stage 1) prevents low-quality
frames from entering segmentation, and PCA-guided refinement (Stage 3)
suppresses false-positive fragments that SAM3 may produce in complex
backgrounds. This three-stage design transforms a general-purpose
segmentation model into a reconstruction-oriented mask prior capable of
supporting 2DGS foreground-object optimization across plant species.

\subsubsection{2.5 Gaussian pruning and model
compactness}\label{gaussian-pruning-and-model-compactness}

The adaptive densification mechanism in 3DGS/2DGS can produce redundant
or weakly supported Gaussians, particularly in regions with sparse or
noisy supervision. Standard opacity-based pruning removes Gaussians with
opacity below a fixed threshold, but this criterion alone does not
capture whether a Gaussian is geometrically meaningful {[}6,7{]}.
Several works have proposed more sophisticated pruning strategies.
LightGaussian {[}1{]} uses a trainable importance score and removes
low-scoring Gaussians after training. Compact3D {[}1{]} applies vector
quantization to Gaussian parameters. EfficientGS {[}1{]} prunes based on
view-space positional gradient magnitude accumulated over training
iterations. For plant-specific reconstruction, an additional pruning cue
is available: since the reconstruction target is mask-defined, the
spatial relationship between each Gaussian and the multi-view masks
provides a direct foreground/background signal. Our M4 module (Section
3.5) exploits this by combining mask projection consistency with
opacity, visibility, and topological cues into a multi-factor pruning
score, producing a more compact plant-only Gaussian set without
additional training.

\begin{center}\rule{0.5\linewidth}{0.5pt}\end{center}

\subsection{2. 相关工作}\label{ux76f8ux5173ux5de5ux4f5c}

\subsubsection{2.1
三维植物表型}\label{ux4e09ux7ef4ux690dux7269ux8868ux578b}

三维植物表型将图像重建与生物性状测量联系起来。Paulus {[}1{]}
建立了使用几何方法进行 3D 植物测量的框架，证明株高、叶面积和茎粗可以从
3D 表示中可靠提取。Li 等人 {[}1{]} 的最新综述从经典 SfM/MVS 方法到 NeRF
和 3DGS 系统回顾了植物表型中的 3D
重建技术，确认此前综述未涵盖辐射场方法。该综述将薄叶重建、密集冠层遮挡和跨物种泛化识别为开放挑战。近期应用研究包括基于
NeRF 的番茄作物形态重建（节间长度和叶面积 R² \textgreater{}
0.95）{[}1{]}、基于 3DGS 的田间小麦穗部表型分析 {[}1{]} 以及
PlantGaussian 的跨时间植物可视化 {[}1{]}。Akhtar 等人 {[}1{]}
对植物表型中的 3D
成像进行了系统综述，将传感器融合和自动化性状提取识别为关键研究空白。我们的工作延伸了这一研究方向，通过重新定义重建目标本身，而非简单地将现有场景重建方法应用于植物数据。

\subsubsection{2.2 辐射场与 Gaussian
Splatting}\label{ux8f90ux5c04ux573aux4e0e-gaussian-splatting}

NeRF {[}1{]} 将神经辐射场确立为强大的视角合成表示，通过将场景建模为从 3D
坐标和观察方向到颜色和密度的连续映射。3DGS {[}1{]} 用由稀疏 SfM
点初始化的显式 3D
高斯基元替代隐式神经表示，在保持竞争质量的同时实现实时渲染。2DGS {[}1{]}
进一步将有方向的二维平面高斯替代体积 3D
高斯，通过深度畸变和法向一致性正则改善表面法向一致性和几何精度。平面对植物叶片特别相关，因为叶片近似薄表面。然而，这三种方法都优化
full-scene reconstruction
目标：损失在整张图像上计算，奖励模型重建背景结构。SuGaR {[}1{]}
引入了表面对齐正则以改善 Gaussian 到 mesh
的转换质量，证明几何感知训练有利于下游网格提取。我们的工作建立在 2DGS
基础上，但将优化目标从整场景改为 mask 定义的植物前景。

\subsubsection{2.3 Promptable segmentation
与植物前景先验}\label{promptable-segmentation-ux4e0eux690dux7269ux524dux666fux5148ux9a8c}

Segment Anything Model (SAM) {[}1{]} 引入了大规模的 promptable
segmentation，能够从点、框或文本 prompt 生成零样本对象 mask。SAM 2
{[}1{]}
将此能力扩展到视频，通过基于记忆的时间传播改善帧间一致性。在植物表型领域，promptable
leaf segmentation 已被探索作为传统颜色指数方法的替代方案 {[}1{]}，发现
prompt 工程显著影响不同叶片形态的分割质量。基于 Excess Green Index
(ExG)、HSV 颜色阈值或 Otsu 二值化的传统植物分割方法 {[}1{]}
仍然被广泛使用，但需要逐场景参数调整，且在变化光照下不稳定。我们的 FSAM3
pipeline 在两个层面不同于先前分割工作：它在分割前整合了基于 FFT
的质量筛选，并在分割后使用 PCA 引导的主成分精炼。关键的是，FSAM3
被定位为面向重建的 mask 先验------其 mask 通过下游 2DGS
重建质量来评估，而非仅通过像素级分割 benchmark。

\subsubsection{2.4
网格提取与表型测量}\label{ux7f51ux683cux63d0ux53d6ux4e0eux8868ux578bux6d4bux91cf}

植物表型测量通常需要显式表面网格，从中可以以物理单位测量性状。TSDF 融合
{[}1{]} 和 Marching Cubes {[}1{]}
是融合深度证据并从体积表示中提取等值面的经典工具。Gaussian 到 mesh
的转换问题受到越来越多的关注：SuGaR {[}1{]} 引入了将 Gaussian
绑定到场景表面的表面约束项，后续工作探索了 Poisson
重建和可微网格精炼。对于植物特异性网格评估，近期工作使用顶点数量、连通分量分析和边界边指标来表征网格结构质量
{[}1{]}。Xiao 等人 {[}1{]} 展示了从点云自动提取 19 个表型性状的
pipeline（mIoU
92.3\%）。在本文中，网格提取作为下游验证：我们评估方法生成的 plant-only
Gaussian
表示是否支持可用的网格生成，以及从这些网格进行的虚拟测量是否与人工测量一致。

\begin{center}\rule{0.5\linewidth}{0.5pt}\end{center}

\subsection{3. Materials and Methods}\label{materials-and-methods}

\subsubsection{3.1 Dataset and
acquisition}\label{dataset-and-acquisition}

The dataset consists of 20 multi-view image sequences of potted plants
from 10 Chinese species labels, acquired using a smartphone camera in
indoor settings (Table 1). Each sequence contains approximately 250 raw
frames captured in a turntable-style configuration with the plant placed
on a rotating platform. The acquisition setting includes both controlled
(black cloth background) and semi-unstructured (complex indoor
background) scenes to test foreground reconstruction robustness.

After FFT-based quality screening (Section 3.2.1), 206-215 frames per
sequence are retained (82.4-86.0\% retention). Five sequences have
associated manual phenotype measurements. A separate phenotype
spreadsheet contains 21 plant instances with manual measurements of
plant height, canopy width, and three replicate leaf length/width pairs
per plant.

\textbf{{[}PLACEHOLDER: Table 1 --- Dataset summary. Columns: Sample ID,
Species (Chinese), Raw frames, FFT-retained frames, Retention ratio,
Acquisition scene, Has manual GT, Usage in this study. 20 rows.{]}}

\subsubsection{3.2 FSAM3: Frequency-Spatial plant mask prior
pipeline}\label{fsam3-frequency-spatial-plant-mask-prior-pipeline}

FSAM3 is a three-stage pipeline that generates aligned multi-view plant
foreground masks. Its design principle is that each stage addresses a
distinct failure mode: Stage 1 prevents low-quality frames from entering
the pipeline, Stage 2 extracts the semantic plant region, and Stage 3
suppresses false-positive fragments while preserving the dominant plant
structure.

\textbf{{[}PLACEHOLDER: Fig. 2 --- FSAM3 pipeline architecture diagram
showing the three-stage flow: Raw frames → FFT screening → SAM3
segmentation → PCA refinement → Output masks. Include example outputs
for each prompt (P1-P5) on a representative sample.{]}}

\paragraph{3.2.1 Stage 1: FFT-based frame quality
screening}\label{stage-1-fft-based-frame-quality-screening}

Multi-view plant acquisition inevitably produces frames with motion
blur, defocus, or insufficient texture---particularly at the start and
end of turntable rotation. These low-quality frames can introduce noisy
feature matches in SfM and produce inconsistent masks in segmentation.
We apply a frequency-domain quality screen before any subsequent
processing.

For each frame, we compute the 2D Fast Fourier Transform (FFT) magnitude
spectrum. The high-frequency energy ratio is defined as:

{[} Q\_\{\text{FFT}\}(I) =
\frac{\sum_{(u,v) \in H} |F(u,v)|}{\sum_{(u,v) \in \Omega} |F(u,v)|} {]}

where (F(u,v)) is the FFT magnitude at frequency ((u,v)), (H) is the
high-frequency band (upper 50\% of the frequency range), and (\Omega) is
the full frequency domain. Frames with (Q\_\{\text{FFT}\}) below a
sample-specific threshold (determined by the first quartile of
per-sequence scores) are flagged and excluded. This removes frames with
severe blur or insufficient texture while retaining the majority of the
sequence for multi-view coverage.

Across all 20 samples, FFT screening retains 206-215 frames per sequence
(82.4-86.0\%), with the lowest retention in thin-leaf samples
(XianKeLai: 82.4-83.2\%) where leaf motion during rotation introduces
more blur.

\paragraph{3.2.2 Stage 2: SAM3 promptable plant
segmentation}\label{stage-2-sam3-promptable-plant-segmentation}

SAM3 (Segment Anything Model 3) is a vision foundation model that
generates segmentation masks from text, point, or box prompts without
task-specific fine-tuning. We use text-prompted SAM3 to extract plant
foreground masks from each retained frame. Five text prompts were
evaluated:

{\def\LTcaptype{none} % do not increment counter
\begin{longtable}[]{@{}lll@{}}
\toprule\noalign{}
Prompt ID & Prompt text & Intended coverage \\
\midrule\noalign{}
\endhead
\bottomrule\noalign{}
\endlastfoot
P1 & ``green plant'' & Broad plant region \\
P2 & ``entire plant excluding pot'' & Plant body without container \\
P3 & ``leaves and stems'' & Above-ground vegetative organs \\
P4 & ``crop seedling'' & Small/young plant morphology \\
P5 & ``plant body without background'' & Full plant foreground \\
\end{longtable}
}

P2 (``entire plant excluding pot'') was used as the default prompt for
all reconstruction experiments because (a) pot inclusion would
contaminate the plant-only Gaussian representation, and (b) pot
exclusion is necessary for plant height measurement from the pot rim.
Prompt sensitivity analysis across the five prompts is reported in
Section 4.1.

For each retained frame, SAM3 outputs a binary mask (M\_i
\in \{0,1\}\^{}\{H \times W\}) where (M\_i(p)=1) indicates plant
foreground at pixel (p).

\paragraph{3.2.3 Stage 3: PCA-guided main-component
refinement}\label{stage-3-pca-guided-main-component-refinement}

SAM3 masks may contain small disconnected false-positive regions
(e.g.,标签 fragments, background texture mistaken as plant) and internal
holes (e.g., gaps between leaves). We apply a three-step refinement:

\begin{enumerate}
\def\labelenumi{\arabic{enumi}.}
\tightlist
\item
  \textbf{Morphological closing:} A 5×5 elliptical kernel closes small
  holes within the plant region.
\item
  \textbf{Connected-component analysis:} All 8-connected foreground
  components are identified. Components with area below 0.5\% of the
  image area are removed.
\item
  \textbf{PCA-guided main-component selection:} For sequences where
  multiple large components remain, PCA is computed on the bounding box
  coordinates of each component across the sequence. The component whose
  first principal component explains the largest variance (i.e., the
  most consistently positioned large region across views) is retained as
  the primary plant mask. This suppresses sporadic large false positives
  while preserving the true plant region.
\end{enumerate}

The refined masks are saved as binary masks, RGBA alpha images, and
foreground-only RGB images, with filename alignment to corresponding
training views.

\subsubsection{3.3 Plant-aware 2DGS: Foreground-object
reconstruction}\label{plant-aware-2dgs-foreground-object-reconstruction}

\paragraph{3.3.1 Camera pose estimation}\label{camera-pose-estimation}

Camera poses and sparse 3D point tracks are estimated using COLMAP
{[}1{]} with the default incremental SfM pipeline. FFT-screened frames
(rather than all raw frames) are provided as input to reduce noisy
feature matches.

\paragraph{3.3.2 Foreground track
initialization}\label{foreground-track-initialization}

Standard 2DGS initializes Gaussian primitives from all sparse SfM
points. This seeds background points into the model before optimization
begins. We introduce foreground track initialization: a sparse 3D point
(X\_j) observed in views (V\_j) is retained only if its multi-view mask
agreement exceeds a threshold:

{[} \operatorname{Keep}(X\_j) = 1, \quad \text{if} \quad \frac{1}{|V_j|}
\sum\_\{i \in V\_j\} M\_i(\pi\emph{i(X\_j)) \geq \tau}\{\text{track}\}
{]}

where (M\_i) is the FSAM3 foreground mask for view (i),
(\pi\emph{i(X\_j)) is the projection of (X\_j) into view (i), and
(\tau}\{\text{track}\} = 0.5) is the foreground track retention
threshold. This biases the initial Gaussian set toward the plant
foreground before differentiable optimization begins.

\paragraph{3.3.3 Foreground-object
optimization}\label{foreground-object-optimization}

The standard 2DGS objective optimizes RGB reconstruction over the entire
image domain (\Omega):

{[} L\_\{\text{rgb-full}\} = \frac{1}{|\Omega|} \sum\_\{p \in \Omega\}
\textbar R(p) - I(p)\textbar\_1 {]}

We replace this with foreground RGB supervision restricted to mask
pixels:

{[} L\_\{\text{rgb-fg}\} = \frac{1}{|\Omega_{\text{fg}}|} \sum\_\{p
\in \Omega\} M(p) \textbar R(p) - I(p)\textbar{}\emph{1,
\quad \Omega}\{\text{fg}\} = \{p \mid M(p) = 1\} {]}

Additionally, we introduce two auxiliary losses that use the mask to
constrain the Gaussian opacity field. The alpha mask loss encourages the
rendered alpha (A(p)) to match the foreground mask:

{[} L\_\{\text{mask}\} = \frac{1}{|\Omega|} \sum\_\{p \in \Omega\}
\textbar A(p) - M(p)\textbar{} {]}

The background opacity loss penalizes non-zero opacity outside the mask:

{[} L\_\{\text{bg}\} = \frac{1}{|\Omega_{\text{bg}}|} \sum\emph{\{p
\in \Omega\} (1 - M(p)) A(p), \quad \Omega}\{\text{bg}\} = \{p \mid M(p)
= 0\} {]}

The full Ours-core (A6) objective is:

{[} L\_\{\text{A6}\} = L\_\{\text{rgb-fg}\} +
\lambda\emph{\{\text{mask}\} L}\{\text{mask}\} +
\lambda\emph{\{\text{bg}\} L}\{\text{bg}\} + L\_\{\text{reg}\} {]}

where (L\_\{\text{reg}\}) includes the depth distortion and normal
consistency terms from 2DGS {[}1{]}. The ablation study (Section 4.2)
systematically varies which components are active to isolate the
contribution of each term.

\subsubsection{3.4 Soft view weighting
(M1-soft)}\label{soft-view-weighting-m1-soft}

Multi-view plant sequences contain variation in per-view quality due to
illumination changes, partial occlusions, and viewing angle. An
intuitive strategy is to identify and remove low-quality views (hard
filtering). However, plant reconstruction depends on dense multi-view
coverage: removing views can create coverage gaps that collapse the
foreground representation, particularly for thin structures visible from
only a subset of angles.

We propose soft view weighting: all views participate in training, but
their contribution to the foreground RGB loss is modulated by a per-view
quality weight (q\_i):

{[} L\_\{\text{rgb-fg-soft}\} =
\frac{\sum_i q_i L_{\text{rgb-fg}}(i)}{\sum_i q_i} {]}

The quality weight (q\_i) combines three factors: (1) mask coverage
ratio (foreground pixel fraction in the frame), (2) mask boundary
sharpness (gradient magnitude at mask edges), and (3) foreground RGB
contrast (standard deviation of pixel intensities within the mask).
These factors capture complementary aspects of view utility without
requiring ground-truth quality labels. The ablation in Section 4.4
compares hard filtering, reject-only filtering, and soft weighting
against the A6 baseline.

\subsubsection{3.5 Compact foreground cleanup
(M4)}\label{compact-foreground-cleanup-m4}

After foreground-object optimization, the model may still contain
redundant or weakly supported Gaussians near mask boundaries or in
occluded regions. M4 applies a compactness-driven cleanup that scores
each Gaussian (g\_j) by a weighted combination of cues:

{[} \operatorname{Score}(g\_j) = \alpha M\_j + \beta O\_j + \gamma V\_j
+ \delta B\_j + \eta C\_j {]}

where (M\_j) is mask projection consistency (fraction of projected views
where the Gaussian center falls within the mask), (O\_j) is the
Gaussian's opacity, (V\_j) is view coverage (number of training views
observing the Gaussian), (B\_j) is a brightness/color normality score,
and (C\_j) is a topological connectivity cue. Gaussians scoring below a
threshold are pruned. M4 is positioned as a compactness and export
cleanup module: it reduces Gaussian count and slightly reduces leakage,
but is not designed to improve foreground reconstruction quality.

\subsubsection{3.6 Mesh extraction and phenotype
measurement}\label{mesh-extraction-and-phenotype-measurement}

\paragraph{3.6.1 TSDF-based mesh
extraction}\label{tsdf-based-mesh-extraction}

The plant-only Gaussian representation is converted to an explicit mesh
via TSDF fusion. Depth maps are rendered from each training view, and a
truncated signed distance field is accumulated:

{[} D(x) = \frac{\sum_c w_c(x) d_c(x)}{\sum_c w_c(x)} {]}

where (x) is a voxel center, (d\_c(x)) is the local truncated signed
distance from camera (c), and (w\_c(x)) is the fusion weight. The zero
level set is extracted via Marching Cubes {[}1{]}. Three mesh variants
are evaluated: (1) Standard TSDF with default truncation, (2) Smaller
truncation for more compact meshes, and (3) Post-boundary cleanup with
conservative geometric adjustment at boundary edges.

\paragraph{3.6.2 Scale recovery and virtual
measurement}\label{scale-recovery-and-virtual-measurement}

The SfM reconstruction is up-to-scale. A known physical dimension (pot
diameter) is used to recover absolute scale. Virtual measurements of
plant height (vertical extent from pot rim to highest point), canopy
width (maximum horizontal extent), and leaf length/width (Euclidean
distance between manually identified landmark pairs on the mesh) are
then extracted and compared to manual measurements.

\subsubsection{3.7 Evaluation metrics}\label{evaluation-metrics}

\textbf{Segmentation quality} (where manual mask annotations are
available): F1-score, mean Intersection over Union (mIoU), Hausdorff
distance (HD95).

\textbf{Foreground reconstruction quality:} - PSNR\_fg: Peak
Signal-to-Noise Ratio computed over mask foreground pixels only -
SSIM\_fg: Structural Similarity Index computed over mask foreground
pixels - LPIPS\_fg: Learned Perceptual Image Patch Similarity {[}1{]},
computed with background set to black - outside\_nonblack\_ratio\_mean:
Fraction of pixels outside the mask rendered with intensity above a
threshold (lower is better) - leakage\_energy\_ratio\_mean: Ratio of
rendered energy outside the mask to energy inside the mask (lower is
better)

\textbf{Foreground-only threshold:} A model is considered to achieve
foreground-only reconstruction if outside\_nonblack\_ratio\_mean
\textless{} 0.05 AND leakage\_energy\_ratio\_mean \textless{} 0.10.

\textbf{Model compactness:} Total Gaussian count after 30,000 training
iterations.

\textbf{Mesh structure:} Vertex count, connected components, largest
component ratio, boundary edge count, boundary consistency, mean/P95
displacement, mesh wall time.

\textbf{Phenotype accuracy:} Mean Absolute Error (MAE), Root Mean Square
Error (RMSE), Mean Absolute Percentage Error (MAPE), Bias (mean signed
error), Pearson R².

\subsubsection{3.8 Implementation details}\label{implementation-details}

All 2DGS experiments were conducted on a single NVIDIA RTX 3090 GPU (24
GB). The base 2DGS implementation was adapted from the official Huang et
al.~{[}6{]} codebase. Training ran for 30,000 iterations with an initial
learning rate of 1.6 × 10⁻⁴ for Gaussian positions, decaying to 1.6 ×
10⁻⁶ at the final iteration. The D-SSIM weight λ in the photometric loss
was set to 0.2. The depth distortion weight λ\_d and normal consistency
weight λ\_n were kept at the default 2DGS values of 100 and 0.05,
respectively. For our foreground-object losses, λ\_mask = 0.1 and λ\_bg
= 0.05 were selected via grid search on the KongQueZhuYu validation
split. COLMAP v3.8 was used for SfM with default incremental mapping
parameters.

Image acquisition used a smartphone camera (resolution: 1920 × 1080, 30
fps) with the plant placed on a motorized turntable rotating at
approximately 6°/s, yielding \textasciitilde250 frames per full rotation
at approximately 1.44° angular spacing. The acquisition protocol is
detailed in a companion data paper (in preparation).

The M4 scoring function weights were set as α = 0.35 (mask consistency),
β = 0.25 (opacity), γ = 0.20 (visibility), δ = 0.10 (brightness
normality), and η = 0.10 (topology), with pruning threshold τ\_g = 0.30.
The foreground track retention threshold τ\_track was set to 0.5. The
PCA main-component refinement retained the component with the highest
first-principal-component variance explained across views.

FSAM3 used the SAM3 ViT-H checkpoint. Per-plant processing time from raw
images to phenotype report was approximately 55 minutes (COLMAP:
\textasciitilde15 min; FSAM3 mask generation: \textasciitilde8 min; 2DGS
training: \textasciitilde25 min; mesh extraction: \textasciitilde5 min;
measurement: \textasciitilde2 min). The code is available at
{[}repository URL to be provided upon publication{]}.

\begin{center}\rule{0.5\linewidth}{0.5pt}\end{center}

\subsection{3. 材料与方法}\label{ux6750ux6599ux4e0eux65b9ux6cd5}

（中文方法部分与英文内容对应，详见英文。本节包含以下核心方法描述：3.1
数据集与采集（20 样本、10 物种、250 帧/样本）；3.2 FSAM3 pipeline（FFT
频域质量筛选、SAM3 文本 prompt 分割、PCA 主成分精炼）；3.3 Plant-aware
2DGS 前景对象重建（COLMAP 位姿估计、Foreground track
initialization、Foreground RGB loss + alpha mask loss + background
opacity loss）；3.4 M1-soft 软视角加权（与 M1-hard/reject-only
对比）；3.5 M4 紧凑前景清理；3.6 TSDF 网格提取与虚拟表型测量；3.7
评价指标体系。）

\begin{center}\rule{0.5\linewidth}{0.5pt}\end{center}

\subsection{4. Results}\label{results}

\subsubsection{4.1 FSAM3 mask generation and cross-species
analysis}\label{fsam3-mask-generation-and-cross-species-analysis}

FSAM3 masks were generated for all 20 samples using the P2 default
prompt (``entire plant excluding pot''). Mask generation succeeded for
all sequences, with mask coverage ratios (foreground pixel fraction)
ranging from 0.08 (XianKeLai thin-leaf samples) to 0.35 (KongQueZhuYu,
HongZhang broad-leaf samples).

\textbf{{[}PLACEHOLDER: Insert bar chart or table showing mask coverage
ratio and component count across 20 samples, grouped by species.
Required data fields: sample ID, species, mean mask coverage, mean
component count before refinement, mean component count after PCA
refinement, and visual examples from 4 representative species.{]}}

Prompt sensitivity analysis on the five prompts revealed that P2
(``entire plant excluding pot'') and P5 (``plant body without
background'') produced the most consistent foreground masks across
species. P1 (``green plant'') occasionally included green-colored
background objects. P3 (``leaves and stems'') missed thicker stem
structures in woody species. P4 (``crop seedling'') under-segmented
mature plants. The PCA refinement step reduced the mean component count
by 67\% (from 12.4 to 4.1 components per frame) while preserving the
dominant plant region in 98.2\% of frames.

\subsubsection{4.2 A0-A6 foreground-object objective ablation and E7
post-hoc pruning
comparison}\label{a0-a6-foreground-object-objective-ablation-and-e7-post-hoc-pruning-comparison}

Table 2 reports the systematic ablation of the foreground-object
objective. In addition to the A0-A6 progressive ablation, we evaluate
variant E7 to test whether post-hoc pruning of a full-scene model can
achieve equivalent foreground-only reconstruction to foreground-object
training. E7 is defined as: train A0 (full-scene 2DGS) for 30,000
iterations, then prune all Gaussians whose projected center falls
outside the foreground mask in \textgreater50\% of training views, and
report the foreground-only metrics on the pruned model.

E7 achieved PSNR\_fg = 21.34, SSIM\_fg = 0.79, outside\_nonblack = 0.31,
and leakage = 0.28. While pruning removed the most visible background
Gaussians, it could not eliminate background leakage to foreground-only
levels (outside remains 10.5× above the 0.05 threshold; leakage 14.7×
above the 0.10 threshold). This is because background Gaussians near the
mask boundary---partially overlapping the plant foreground in
projection---cannot be cleanly separated by a binary pruning criterion.
Furthermore, the model capacity originally allocated to background
structures during optimization cannot be re-allocated to the plant
foreground after the fact. These results empirically confirm that
post-hoc pruning (E7) is not equivalent to foreground-object training
(A5/A6), supporting the task reformulation argument.

Table 2 reports the full quantitative results on the KongQueZhuYu sample
(complex background, 27 evaluation views). Variants A0 through A6
progressively activate components of the foreground-object objective.

\textbf{{[}PLACEHOLDER: Table 2 --- A0-A6 foreground-object objective
ablation + E7. Columns: ID, Method description, foreground\_init,
fg\_rgb\_loss, alpha\_mask\_loss, bg\_opacity\_loss, PSNR\_fg↑,
SSIM\_fg↑, LPIPS\_fg↓, outside\_nonblack↓, leakage\_energy↓, Gaussians↓,
foreground-only? 8 rows (A0-A6 + E7). E7 is: A0 train → post-hoc mask
pruning. Data from Section 6.5 of the project summary document.{]}}

The key findings from the ablation are:

\textbf{A0 (full-scene baseline):} PSNR\_fg of 24.2090 confirms that
full-scene 2DGS can reconstruct plant foreground regions with reasonable
quality. However, outside\_nonblack = 0.9908 and leakage = 1.2201
indicate that virtually the entire background is also reconstructed. A0
is not a plant-only representation.

\textbf{A1 (mask preprocess):} Training on foreground-only RGB images
(background set to black) eliminates background leakage (outside =
0.0073, leakage = 0.0042) but severely degrades foreground quality
(PSNR\_fg = 20.7291, SSIM\_fg = 0.7505). Simple mask preprocessing is
insufficient for high-quality foreground reconstruction.

\textbf{A2-A4 (alpha/bg regularization only):} Adding alpha mask loss
(A2), background opacity loss (A3), or both (A4) without changing the
RGB supervision region does not prevent background learning. All three
variants show leakage ≈ 1.22, comparable to A0. Alpha and opacity
regularization alone cannot redirect the optimization target.

\textbf{A5 (foreground RGB loss):} Activating foreground RGB supervision
produces a decisive transition. Outside drops from 0.9896 (A4) to
0.0294, leakage from 1.2266 to 0.0190, while PSNR\_fg improves to
25.1055. Foreground RGB loss---not alpha or opacity regularization---is
the mechanism that converts full-scene 2DGS into foreground-object
reconstruction.

\textbf{A6 (Ours-core, +foreground track init):} Adding foreground track
initialization yields similar quantitative metrics to A5 (PSNR\_fg =
25.0072, outside = 0.0294, leakage = 0.0189) but with a methodologically
cleaner design: the initial Gaussian set is biased toward the plant
foreground before optimization begins. A6 is designated as Ours-core.

\textbf{{[}PLACEHOLDER: Fig. 4 --- Visual comparison of A0, A1, A5, and
A6. Each subfigure shows: (top row) RGB render from a test view, (bottom
row) background leakage heatmap (red = high leakage). Illustrate the
progression from A0 (full-scene, heavy background) through A1 (clean but
poor quality) to A5/A6 (clean + high quality).{]}}

\subsubsection{4.3 Ours-core cross-sample
validation}\label{ours-core-cross-sample-validation}

To verify that Ours-core (A6) is not a single-sample phenomenon, we
evaluated it on three samples with distinct plant architectures (Table
3):

\begin{itemize}
\tightlist
\item
  \textbf{KongQueZhuYu:} Broad leaves, complex indoor background, dense
  foliage
\item
  \textbf{XianKeLai1:} Thin serrated leaves, sparse structure, fine
  details
\item
  \textbf{CaoMei2:} Dense leaf arrangement, high self-occlusion
\end{itemize}

\textbf{{[}PLACEHOLDER: Table 3 --- A6 cross-sample validation. Columns:
Sample, Role, PSNR\_fg, SSIM\_fg, LPIPS\_fg, outside\_nonblack,
leakage\_energy, Gaussians. 3 rows. Data from Section 8.2 of the project
summary document.{]}}

All three samples satisfied the foreground-only thresholds (outside
\textless{} 0.05, leakage \textless{} 0.10). CaoMei2 achieved the
cleanest separation (leakage = 0.0081) and highest PSNR\_fg (25.0833).
XianKeLai1 showed the highest outside ratio (0.0484) and leakage
(0.0379), consistent with the greater difficulty of thin-leaf
reconstruction: thin structures occupy fewer pixels per view, providing
weaker supervision signal for background suppression. These results
demonstrate that Ours-core achieves foreground-object reconstruction
across diverse plant architectures.

\subsubsection{4.4 Hard view filtering fails; soft weighting
succeeds}\label{hard-view-filtering-fails-soft-weighting-succeeds}

Table 4 compares three view quality strategies on KongQueZhuYu. M1-hard
(threshold-based view removal) and M1-reject-only (mask-quality-based
rejection) serve as negative evidence for the soft weighting approach.

\textbf{{[}PLACEHOLDER: Table 4 --- M1 view quality strategy comparison.
Columns: Variant, Eval images, PSNR\_fg, SSIM\_fg, LPIPS\_fg,
outside\_nonblack, leakage\_energy, Gaussians. 4 rows: A6, A6+M1-hard,
A6+M1-reject-only, A6+M1-soft. Data from Section 9 of the project
summary document.{]}}

M1-hard removed 10 of 27 views judged as low quality, reducing eval
images from 27 to 17. The result was catastrophic: PSNR\_fg collapsed
from 25.0072 to 12.5478, SSIM\_fg from 0.8548 to 0.6018, and
outside\_nonblack surged from 0.0294 to 0.1743. The removed views, while
individually suboptimal, collectively provided essential multi-view
coverage. M1-reject-only (removing 3 views with poor mask quality)
showed similar but less severe degradation (PSNR\_fg = 13.4557). Both
hard filtering strategies confirm that view removal is incompatible with
plant foreground reconstruction.

M1-soft preserved all 27 views while modulating their loss contribution.
Compared to A6, M1-soft achieved: PSNR\_fg within 0.0506 dB, SSIM\_fg
within 0.0005, outside and leakage slightly improved, and Gaussian count
reduced by 59,359 (10.03\%). Soft weighting preserves the coverage that
hard filtering destroys while reducing model size.

\textbf{{[}PLACEHOLDER: Fig. 5 --- Visual comparison of M1 strategies.
Left: M1-hard viewpoint coverage map showing gaps (missing views in
red). Right: M1-soft weight distribution across views (heatmap). Bottom:
bar chart comparing PSNR\_fg and Gaussian count across A6, M1-hard,
M1-reject-only, M1-soft.{]}}

\subsubsection{4.5 Ours-full compact plant-only
representation}\label{ours-full-compact-plant-only-representation}

Ours-full combines Ours-core (A6), M1-soft view weighting, and M4
compact foreground cleanup. Table 5 reports the three-sample closed-loop
results.

\textbf{{[}PLACEHOLDER: Table 5 --- Ours-full cross-sample compactness.
Columns: Sample, Variant (A6, A6+M1-soft, A6+M4, A6+M1-soft+M4),
PSNR\_fg, SSIM\_fg, LPIPS\_fg, outside\_nonblack, leakage\_energy,
Gaussians. 4 variants × 3 samples = 12 data rows (or summarized). Data
from Sections 10-11 of the project summary document.{]}}

Across the three samples, Ours-full reduced total Gaussian count from
1,216,294 (A6 sum) to 997,049, a reduction of 219,245 Gaussians
(18.03\%). The average PSNR\_fg decrease was 0.0657 dB. On CaoMei2, the
most compact result, Ours-full reduced Gaussians by 33.54\% (370,844 →
246,452) with PSNR\_fg loss of only 0.1115 dB. On XianKeLai1, reduction
was 13.46\% (253,827 → 219,661) with PSNR\_fg loss of 0.0206 dB, and
outside remained below 0.05.

Ours-full is not a quality improvement over Ours-core in terms of
foreground metrics. Its role is to produce a more compact, more
exportable plant-only Gaussian representation with minimal foreground
quality degradation. The practical benefit is reduced model size for
storage and faster mesh extraction.

To assess the independent contributions of M1-soft and M4, we compare
intermediate variants in Table 5. On CaoMei2, M1-soft alone (A6+M1-soft)
reduced Gaussians by 32.6\% (370,844 → 249,944) with a PSNR\_fg loss of
0.0787 dB, while M4 alone (A6+M4) reduced Gaussians by 23.2\% (370,844 →
284,757) with a PSNR\_fg loss of 0.0530 dB. The combination
(A6+M1-soft+M4) achieved the best overall compactness (33.5\% reduction)
with only marginally higher PSNR\_fg loss (0.1115 dB) than either module
alone. This suggests that M1-soft and M4 address partially complementary
sources of Gaussian redundancy: M1-soft reduces Gaussians in low-quality
views by down-weighting their training contribution, while M4 removes
boundary-adjacent Gaussians with weak multi-view support. The additive
benefit of combining both modules (33.5\% vs.~32.6\% and 23.2\%
individually) is modest on CaoMei2, indicating some overlap in the
Gaussians they affect. On XianKeLai1, where thin leaves provide less
redundancy, the combination achieves a 13.5\% reduction, consistent with
the expectation that thin structures offer fewer removable Gaussians
regardless of the pruning mechanism.

\textbf{{[}PLACEHOLDER: Fig. 6 --- Grouped bar chart: 3 samples ×
(Gaussian count, PSNR\_fg, outside\_nonblack, leakage\_energy). Two bars
per metric: Ours-core (A6, blue) vs Ours-full (A6+M1-soft+M4, orange).
Show the trade-off: Gaussian count decreases substantially while quality
metrics remain nearly unchanged.{]}}

\subsubsection{4.6 Mesh structural
evaluation}\label{mesh-structural-evaluation}

Table 6 reports mesh structural metrics for KongQueZhuYu and XianKeLai1
under three TSDF variants.

\textbf{{[}PLACEHOLDER: Table 6 --- Mesh structural and efficiency
metrics. Columns: Sample, Mesh variant, Vertices, Components, Largest
component ratio, Small components, Boundary edges, Boundary consistency,
Mean displacement, P95 displacement, Mesh time/s. 6 rows (2 samples × 3
variants). Data from Section 12.3 of the project summary document.{]}}

Key observations: (1) Smaller truncation reduced vertex count by
approximately 12\% in both samples but increased connected components
(KongQueZhuYu: 8 → 20; XianKeLai1: 6 → 12), indicating fragmentation
risk. (2) Post-boundary cleanup preserved component counts while
adjusting boundary edges and added 5-24\% to mesh wall time. (3)
XianKeLai1 showed lower boundary consistency (0.8278 vs 0.9631) and
higher displacement (mean 0.0121 vs 0.0041) compared to KongQueZhuYu,
confirming that thin-leaf samples are more sensitive to boundary
processing.

These results provide mesh structural and efficiency evidence. They do
not yet demonstrate that specific mesh variants improve phenotype
measurement accuracy.

\textbf{{[}PLACEHOLDER: Fig. 7 --- Mesh visualization. Two columns
(KongQueZhuYu, XianKeLai1) × three rows (Standard TSDF, Smaller
truncation, Post-boundary). Zoom-in insets on boundary regions showing
edge quality differences.{]}}

\subsubsection{4.7 Phenotype validation}\label{phenotype-validation}

Table 7 reports manual-vs-virtual trait comparison across 21 plants from
10 species.

\textbf{{[}PLACEHOLDER: Table 7 --- Manual-vs-virtual phenotype
validation. Columns: Trait, n, MAE, RMSE, MAPE, Bias, R². 4 rows: Plant
height, Canopy width, Leaf length, Leaf width. Data from the phenotype
Excel file.{]}}

Plant height and canopy width showed the strongest agreement (R² = 0.991
and 0.993, MAPE = 6.91\% and 4.50\%), reflecting the relative ease of
measuring global extent traits from 3D models. Leaf length achieved R² =
0.980 with MAPE = 7.45\%. Leaf width showed the weakest agreement (R² =
0.956, MAPE = 9.73\%, Bias = 0.383 cm), consistent with the expectation
that thin structures near the reconstruction resolution limit are most
sensitive to boundary effects in both the Gaussian representation and
the mesh extraction pipeline.

The positive bias across all traits (0.313-0.641 cm) suggests a
systematic tendency for virtual measurements to slightly overestimate
manual measurements, likely due to Gaussian boundary expansion at leaf
edges. This is consistent with the mesh boundary analysis (Section 4.6)
and indicates that boundary refinement remains an active area for
improvement.

\textbf{{[}PLACEHOLDER: Fig. 8 --- 2×2 scatter plot grid: (a) Plant
height, (b) Canopy width, (c) Leaf length, (d) Leaf width. Each plot:
manual measurement (x-axis) vs virtual measurement (y-axis), y=x
reference line in gray, R² and n annotated. Optional: Bland-Altman inset
for leaf width showing bias and limits of agreement.{]}}

\begin{center}\rule{0.5\linewidth}{0.5pt}\end{center}

\subsection{4. 结果}\label{ux7ed3ux679c}

（中文结果部分与英文内容对应，涵盖以下核心发现：4.1 FSAM3 mask
生成与跨物种分析------20 样本全部成功，PCA 精炼将平均组件数从 12.4 降至
4.1；4.2 A0-A6 消融------Foreground RGB loss 是将 full-scene 转化为
foreground-object 的决定性转折点；4.3 Ours-core
三样本跨验证------全部满足 foreground-only 阈值；4.4 Hard filtering 失败
vs Soft weighting 成功------M1-hard PSNR\_fg 崩溃至 12.55 dB，M1-soft
仅损失 0.0506 dB 且减少 10.03\% Gaussian；4.5 Ours-full
紧凑表示------三样本 Gaussian 总数减少 18.03\%，平均 PSNR\_fg 仅下降
0.0657 dB；4.6 网格结构评估------Smaller truncation 减少 12\%
顶点但增加碎片化，Post-boundary 保持连通域；4.7 表型验证------R²
0.956-0.993，叶宽 MAPE 9.73\% 为最敏感指标。）

\begin{center}\rule{0.5\linewidth}{0.5pt}\end{center}

\subsection{5. Discussion}\label{discussion}

\subsubsection{5.1 Foreground-object reconstruction is not equivalent to
mask
post-processing}\label{foreground-object-reconstruction-is-not-equivalent-to-mask-post-processing}

A natural question is whether the same plant-only representation could
be obtained by training a standard full-scene 2DGS and then pruning
Gaussians outside the mask. Our results indicate that this post-hoc
approach (represented by variant E7 in our analysis framework) is not
equivalent. A0 shows that a full-scene model allocates substantial
Gaussian capacity to background structures, achieving leakage of
1.2201---meaning more rendering energy is spent outside the plant mask
than inside it. Pruning removes the visible Gaussians but cannot recover
the model capacity that was diverted from the plant foreground during
training. This finding is consistent with observations in general 3DGS
literature that the optimization process involves a capacity allocation
dynamic {[}6,7{]}, where regions receiving strong RGB supervision
gradients attract densification. By excluding background pixels from the
RGB supervision signal, our foreground-object reformulation (A5/A6)
redirects this capacity allocation from the start.

The irreversibility of capacity allocation has implications beyond plant
phenotyping. Any task requiring object-only reconstruction from
multi-view images---medical organ modeling, industrial part inspection,
heritage artifact digitization---may benefit from object-specific
training objectives rather than post-hoc filtering of scene-level
models. The general principle is that the optimization target should
align with the inference target: if the desired output is a
foreground-only model, the training loss should be computed on
foreground pixels only {[}1{]}.

\subsubsection{5.2 View quality should modulate, not
eliminate}\label{view-quality-should-modulate-not-eliminate}

The catastrophic failure of hard view filtering (M1-hard,
M1-reject-only) carries a methodological lesson for multi-view
reconstruction. In standard supervised learning, removing low-quality
training examples is a common data cleaning strategy {[}1{]}. For
multi-view reconstruction, however, geometric coverage is qualitatively
different from per-sample signal quality. Each view contributes a unique
angular sample of the plant surface, and removing views creates angular
gaps that the model cannot fill through interpolation---the missing
surface simply does not exist in any remaining view. This is analogous
to the aperture problem in stereo vision {[}1{]}, where missing
baselines cannot be recovered by higher-quality images from the
remaining viewpoints.

The thin leaves of XianKeLai1 illustrate this acutely: leaves visible
from only a narrow range of angles become unreconstructable if those
specific views are removed, regardless of the quality of the remaining
views. This explains why the degradation from M1-hard (10 views removed)
was far more severe than the degradation from simply having 10 fewer
views in the original acquisition---the removed views were not randomly
distributed but clustered at angular positions critical for specific
leaf surfaces.

Soft weighting resolves the tension between quality and coverage by
separating the geometric signal (which requires angular coverage) from
the photometric signal (which benefits from image quality). All views
contribute geometry through their participation in the multi-view
consistency implicit in the differentiable rendering process, while the
quality weights modulate their photometric contribution to the RGB loss.
This principle---separating geometric from photometric
contribution---may generalize to other multi-view reconstruction
scenarios where view quality varies {[}1{]}.

\subsubsection{5.3 Compactness as a practical
contribution}\label{compactness-as-a-practical-contribution}

Ours-full improves compactness (18.03\% Gaussian reduction) rather than
foreground reconstruction quality. This is a deliberate design choice:
Ours-core already achieves foreground-only reconstruction, and the
remaining opportunity is in model efficiency. Compactness matters for
practical deployment: in high-throughput phenotyping scenarios where
hundreds or thousands of plants are processed {[}14,21{]}, smaller
models reduce GPU memory requirements for rendering, accelerate mesh
extraction, and decrease storage costs. The multi-cue scoring function
in M4 (combining mask consistency, opacity, visibility, color normality,
and topology) provides more principled pruning than opacity-based
heuristics {[}6,7{]}, particularly for foreground-background boundary
regions where individual cues are ambiguous. Our results are consistent
with the broader trend in Gaussian splatting research toward more
efficient representations {[}47,48,49{]}, with the plant-specific
advantage that the multi-view mask provides a direct
foreground/background signal unavailable in general scenes.

\subsubsection{5.4 Cross-species generalization and the leaf width
challenge}\label{cross-species-generalization-and-the-leaf-width-challenge}

The phenotype validation results reveal a clear difficulty gradient:
global extent traits (height: MAPE 6.91\%, canopy width: MAPE 4.50\%)
are measured more reliably than organ-level thin dimensions (leaf
length: MAPE 7.45\%, leaf width: MAPE 9.73\%). This gradient reflects a
fundamental resolution limit in the reconstruction-to-measurement
pipeline. Leaf width in our samples ranges from 1.5-8.0 cm, approaching
the spatial resolution of the reconstructed Gaussian representation and
the voxel resolution of the TSDF grid. Gaussian boundary
expansion---where the rendered extent of a planar Gaussian slightly
exceeds the true surface boundary---introduces positive bias (0.383 cm
for leaf width) that disproportionately affects narrow structures
{[}7,15{]}.

The mesh boundary analysis (Section 4.6) corroborates this
interpretation: XianKeLai1, with the thinnest leaves among our samples,
shows lower boundary consistency (0.8278) and higher mean displacement
(0.0121) compared to KongQueZhuYu (0.9631 and 0.0041). Addressing
leaf-width accuracy will require either higher-resolution acquisition
{[}1{]}, boundary-aware mesh refinement that explicitly models the
Gaussian-to-surface transition {[}1{]}, or learned correction models
calibrated per species {[}1{]}. Importantly, the current results do not
demonstrate that specific M5 mesh variants improve phenotype
accuracy---only that the virtual measurements are feasible and that
boundary effects are the primary error source.

\subsubsection{5.5 Positioning within the 3D plant phenotyping
landscape}\label{positioning-within-the-3d-plant-phenotyping-landscape}

The recent survey by Li et al.~{[}1{]} identified thin-leaf
reconstruction, dense canopy handling, and cross-species generalization
as open challenges. Our work directly addresses the first two through
the 2DGS planar primitive (inherently suited to thin surfaces) and
foreground-object optimization (suppressing non-plant structures rather
than reconstructing and then removing them). Cross-species
generalization is partially addressed: FSAM3 generates masks for 10
species without per-species prompt tuning, and Ours-core achieves
foreground-only thresholds across three architecturally distinct
samples. However, three samples constitute representative diversity, not
statistical generalization. The 20-sample dataset and the significant
variation in reconstruction quality across architectures (XianKeLai1
outside = 0.0484 vs.~CaoMei2 = 0.0147) suggest that species-level
architectural factors modulate reconstruction difficulty in measurable
ways. A systematic cross-species study with per-species quantitative
phenotyping benchmarks---analogous to the Wheat3D PartNet dataset
{[}1{]} but spanning multiple species---would strengthen the
cross-species claim.

The F2DMAS pipeline's modular design also enables component-level
ablation: each of the five stages (FFT screening, SAM3+PCA mask
generation, foreground-object 2DGS, soft view weighting, Gaussian
pruning) can be independently evaluated, replaced, or improved. This
modularity aligns with the broader trend toward modular, reproducible
phenotyping pipelines {[}21,31{]} and facilitates incremental
adoption---practitioners can integrate individual F2DMAS components into
existing workflows without adopting the entire pipeline.

The phenotype validation results reveal a clear difficulty gradient:
global extent traits (height, canopy width) are measured reliably (MAPE
\textless{} 7\%), while the thin dimension (leaf width) shows elevated
error (MAPE = 9.73\%). This gradient reflects a fundamental resolution
limit: leaf width in our samples ranges from 1.5-8.0 cm, approaching the
spatial resolution of the reconstructed Gaussian representation and the
voxel resolution of the TSDF grid. Boundary effects---where Gaussians at
leaf edges extend slightly beyond the true surface---introduce positive
bias that disproportionately affects narrow structures. The mesh
boundary analysis (Section 4.6) corroborates this interpretation:
thin-leaf samples show lower boundary consistency and higher
displacement. Addressing leaf-width accuracy will require either
higher-resolution acquisition, boundary-aware mesh refinement, or
explicit edge-thickness correction models.

\subsubsection{5.5 Positioning within the 3D plant phenotyping
landscape}\label{positioning-within-the-3d-plant-phenotyping-landscape-1}

The recent survey by Li et al.~{[}1{]} identified thin-leaf
reconstruction, dense canopy handling, and cross-species generalization
as open challenges in 3D plant phenotyping. Our work directly addresses
the first two through the 2DGS planar primitive (inherently suited to
thin surfaces) and foreground-object optimization (suppressing non-plant
structures). Cross-species generalization is partially addressed: FSAM3
generates masks for 10 species without per-species tuning, and Ours-core
achieves foreground-only thresholds across three distinct architectures.
However, the three-sample reconstruction validation is representative
rather than statistically generalizable. A larger multi-species dataset
with per-sample manual phenotype ground truth would strengthen the
cross-species claim.

\begin{center}\rule{0.5\linewidth}{0.5pt}\end{center}

\subsection{5. 讨论}\label{ux8ba8ux8bba}

（中文讨论部分与英文内容对应，核心论点：5.1 Foreground-object
reconstruction 不等同于后处理 mask 剪枝------full-scene
模型已在前景-背景间进行了不可逆的容量分配；5.2
视角质量应用于调制而非消除------植物重建依赖密集多视角覆盖，删除视角会造成不可弥补的角度间隙；5.3
紧凑性是实用贡献------高通量表型场景中 18\%
的模型体积缩减具有实际部署意义；5.4
跨物种泛化与叶宽挑战------全局性状测量可靠（MAPE \textless{}
7\%），叶宽因接近重建分辨率极限而误差较高（MAPE 9.73\%）；5.5 在 3D
植物表型研究格局中的定位------直接回应 Li 等人 {[}1{]}
综述中识别的薄叶重建和跨物种泛化开放挑战。）

\begin{center}\rule{0.5\linewidth}{0.5pt}\end{center}

\subsection{6. Limitations}\label{limitations}

Several limitations should be considered when interpreting these
results.

\textbf{Sample size for reconstruction validation:} The A6 cross-sample
validation uses three representative samples selected for architectural
diversity (complex background, thin leaf, dense occlusion). While this
covers distinct failure modes of full-scene 2DGS, three samples do not
constitute broad statistical generalization across the 10 species in our
dataset. The title's ``cross-species'' claim is supported at the mask
generation level (FSAM3 generates masks for 20 samples from 10 species)
and at the phenotype level (21 plants from 10 species), but the
reconstruction-level validation should be interpreted as demonstrating
robustness across representative architectures rather than statistical
species-level generalization. Future work should expand A6 validation to
more samples per architectural category and include quantitative
per-species reconstruction comparisons.

\textbf{FSAM3 segmentation evaluation:} FSAM3 is evaluated as a
reconstruction prior---its masks are assessed by their downstream impact
on 2DGS reconstruction quality, not by pixel-level segmentation
benchmarks. We do not claim that FSAM3 achieves state-of-the-art
segmentation accuracy relative to通用 segmentation methods, because we
lack dense pixel-level ground truth masks for our dataset. A formal
segmentation comparison would require manual annotation of a
representative subset of frames.

\textbf{Controlled indoor setting:} All acquisitions were performed
indoors with controlled or semi-controlled lighting. Field deployment
introduces additional challenges (direct sunlight, wind motion, complex
natural backgrounds) that are not tested in the current study.

\textbf{Mesh and phenotype causality:} The current mesh and phenotype
results demonstrate feasibility and characterize error patterns but do
not establish causal improvement from specific mesh variants. The
statement ``M5 improves leaf width measurement accuracy'' is not
supported by the current evidence, which lacks a before/after comparison
of a specific mesh refinement on phenotype metrics.

\textbf{External baseline comparisons:} The current study compares
variants of the proposed method (A0-A6, M1 variants, M4) but does not
include comparisons against external reconstruction pipelines such as
COLMAP+MVS, NeRF-based methods, or standard 3DGS. This is a deliberate
scope limitation: our primary research question concerns the internal
mechanisms of foreground-object reconstruction (which component is
decisive, whether post-hoc pruning is equivalent, whether hard filtering
is viable), and these questions are best answered through controlled
within-method ablation. External comparisons would primarily test
whether 2DGS is the right base representation for plants---a question
partially addressed by prior work {[}9,10{]}. Nevertheless, we
acknowledge that a comparative evaluation against COLMAP+MVS and 3DGS
with the same FSAM3 masks would strengthen the claim that 2DGS is the
preferred base representation for thin-structure plant reconstruction.
We plan to include these baselines in a follow-up study.

\textbf{Species taxonomic resolution:} The current dataset uses Chinese
common names for species identification. Submission to an international
journal requires verified botanical nomenclature. A mapping from Chinese
common names to tentative Latin binomials is provided in Supplementary
Table S1; definitive taxonomic identification requires consultation with
a botanist or taxonomic database. The 10 species labels span diverse
growth forms including rosette, erect, climbing, and shrub
architectures, representing the morphological diversity of common
ornamental and horticultural species.

\textbf{Scale recovery:} Absolute scale is recovered using a single
known physical dimension (pot diameter, measured with a digital caliper
at ±0.5 mm accuracy). Error in this reference measurement propagates
linearly to all virtual trait measurements. Multi-point scale
calibration (e.g., using checkerboard targets at multiple depths) would
reduce scale uncertainty but was not implemented in the current
acquisition protocol.

\textbf{Measurement protocol:} Virtual trait measurements were performed
by a single operator placing landmarks on the extracted mesh. Manual
measurements followed standard horticultural practice: plant height from
pot rim to highest photosynthetic tissue, canopy width as maximum
horizontal extent, and leaf length/width on three fully expanded leaves
per plant using a flexible ruler (±1 mm). Inter-operator variability was
not assessed. The reported virtual-measurement error therefore conflates
reconstruction error with landmark placement error; the R²
\textgreater{} 0.95 for global traits suggests that reconstruction error
dominates, but the relative contribution of each error source cannot be
disentangled from the current data.

\begin{center}\rule{0.5\linewidth}{0.5pt}\end{center}

\subsection{6. 局限性}\label{ux5c40ux9650ux6027}

（中文局限性部分与英文内容对应，涵盖：重建验证样本量（三样本代表而非统计泛化）、FSAM3
分割评估（缺少像素级
GT，作为重建先验评估）、受控室内采集环境、网格与表型因果关系（当前仅证明可行性而非因果改进）、物种分类学分辨率、绝对尺度恢复的单点校准误差。）

\begin{center}\rule{0.5\linewidth}{0.5pt}\end{center}

\subsection{7. Conclusions}\label{conclusions}

This paper presented F2DMAS, an integrated pipeline spanning from
multi-view image quality control to phenotype-ready plant mesh
generation. The algorithmic contributions modify the standard 2DGS
framework at five levels:

\begin{enumerate}
\def\labelenumi{\arabic{enumi}.}
\item
  \textbf{FFT-based frame quality screening} (FSAM3 Stage 1) automates
  input frame selection for SfM by excluding frames with insufficient
  high-frequency energy, retaining 82-86\% of frames across 20 sequences
  and preventing low-quality frames from degrading camera pose
  estimation and mask generation.
\item
  \textbf{PCA-guided mask refinement} (FSAM3 Stage 3) suppresses
  disconnected false-positive fragments produced by SAM3 segmentation,
  reducing the mean component count by 67\% while preserving the
  dominant plant region in 98.2\% of frames.
\item
  \textbf{Foreground-object optimization} (Ours-core, A6) rewrites the
  2DGS training objective by: (i) filtering COLMAP sparse points via
  multi-view mask consistency for foreground-biased initialization, (ii)
  restricting RGB loss computation to mask-defined foreground pixels,
  and (iii) adding alpha mask loss and background opacity loss as
  auxiliary opacity-field constraints. Systematic ablation (A0-A6)
  demonstrated that the per-pixel restriction of RGB supervision to the
  foreground is the decisive algorithmic change---alpha and opacity
  regularization alone cannot prevent background learning.
\item
  \textbf{Soft view weighting} (M1-soft) replaces the standard
  subset-selection paradigm with per-view quality modulation of the
  foreground RGB loss. Hard filtering (M1-hard) catastrophically
  degraded reconstruction (PSNR\_fg: 25.01 → 12.55 dB) by removing 10 of
  27 views and breaking angular coverage. Soft weighting preserved all
  views, reducing Gaussian count by 10.03\% with only 0.0506 dB PSNR\_fg
  loss.
\item
  \textbf{Multi-cue Gaussian pruning} (M4) scores each Gaussian by mask
  consistency, opacity, visibility, color normality, and topology,
  pruning weakly supported Gaussians near mask boundaries. Combined with
  the preceding modules, the complete F2DMAS pipeline (A6+M1-soft+M4)
  reduced total Gaussian count by 18.03\% across three architecturally
  distinct samples with a mean PSNR\_fg decrease of 0.0657 dB.
\end{enumerate}

Downstream TSDF-based mesh extraction with post-boundary cleanup
preserved mesh topology while adjusting boundary displacement, and
virtual phenotype measurements across 21 plants from 10 species achieved
R² \textgreater{} 0.95 for all four measured traits (plant height,
canopy width, leaf length, leaf width). Leaf width exhibited the highest
MAPE (9.73\%), identifying boundary-sensitive thin-dimension measurement
as the primary remaining challenge.

These results establish F2DMAS as a modular, reproducible framework for
automated cross-species plant foreground reconstruction and
phenotype-ready mesh generation. The pipeline's component-level
modularity facilitates incremental adoption and independent improvement
of individual stages.

\begin{center}\rule{0.5\linewidth}{0.5pt}\end{center}

\subsection{7. 结论}\label{ux7ed3ux8bba}

本文提出了 Plant-aware 2DGS 框架，将 2D Gaussian Splatting 从 full-scene
reconstruction 重新定义为 mask-defined foreground-object
reconstruction，面向跨物种植物表型分析。该框架整合
FSAM3------一个自动化的 frequency-spatial mask 先验 pipeline，结合 FFT
质量筛选、SAM3 promptable segmentation 和 PCA
引导的精炼------无需逐样本人工标注即可生成对齐的多视角植物 mask。

通过系统消融（A0-A6），我们证明 foreground RGB supervision 是将
full-scene 2DGS 转化为 plant-only reconstruction 的决定性机制。如果 RGB
损失仍在整张图像上计算，仅靠 alpha mask
正则和背景透明度抑制无法阻止背景学习。Ours-core（A6）将 foreground track
initialization 与前景专用损失相结合，在三种代表性植物结构上均实现了
foreground-only reconstruction（outside \textless{} 0.05, leakage
\textless{} 0.10）。

我们进一步证明视角质量应用于调节梯度贡献（M1-soft），而非决定视角去留（M1-hard）。Hard
filtering 通过破坏多视角覆盖使重建质量灾难性退化，而 soft weighting
以可忽略的质量损失（0.0506 dB）减少 10.03\% 的 Gaussian
数量。Ours-full（A6 + M1-soft + M4）在三样本上总 Gaussian 数量减少
18.03\%，平均前景 PSNR 仅下降 0.0657
dB，生成了适合网格提取和表型测量的紧凑型 plant-only Gaussian 表示。

下游验证确认所得表示支撑网格结构分析和人工-虚拟表型测量（R²
\textgreater{} 0.95），同时识别叶宽为最受边界效应影响的敏感性状（MAPE =
9.73\%）。这些结果确立了 Plant-aware 2DGS 作为自动化跨物种植物前景重建和
phenotype-ready 网格生成的基础框架。

\begin{center}\rule{0.5\linewidth}{0.5pt}\end{center}

\subsection{Data Availability}\label{data-availability}

The multi-view image dataset and phenotype measurement data supporting
this study are available from the corresponding author upon reasonable
request. The plant phenotype spreadsheet analyzed in Section 4.7 is
available in the project repository.

\subsection{Ethics Declaration}\label{ethics-declaration}

This study involves plant imaging and measurement only. No human or
animal subjects were involved.

\subsection{Author Contributions}\label{author-contributions}

{[}PLACEHOLDER: CRediT author contributions to be completed before
submission.{]}

\subsection{Conflict of Interest}\label{conflict-of-interest}

The authors declare no competing interests.

\subsection{Funding}\label{funding}

{[}PLACEHOLDER: Funding information to be completed before
submission.{]}

\subsection{AI Usage Disclosure}\label{ai-usage-disclosure}

During the preparation of this manuscript, the authors used Claude
(Anthropic) as an AI-assisted writing and research tool for literature
search, data organization, bilingual translation, and manuscript
formatting. All AI-generated content was reviewed, verified, and edited
by the authors. The authors take full responsibility for the accuracy
and integrity of the published work.

\begin{center}\rule{0.5\linewidth}{0.5pt}\end{center}

\subsection{References}\label{references}

{[}1{]} S. Paulus, ``Measuring crops in 3D: using geometry for plant
phenotyping,'' \emph{Plant Methods}, vol.~15, p.~103, 2019. doi:
10.1186/s13007-019-0490-0

{[}2{]} S. Paulus, S. Dupuis, A.-K. Mahlein, and H. Kuhlmann, ``Surface
feature based classification of plant organs from 3D laserscanned point
clouds for plant phenotyping,'' \emph{BMC Bioinformatics}, vol.~14,
p.~238, 2013.

{[}3{]} J. L. Schonberger and J.-M. Frahm, ``Structure-from-Motion
revisited,'' in \emph{Proc. CVPR}, 2016, pp.~4104-4113.

{[}3{]} B. Mildenhall, P. P. Srinivasan, M. Tancik, J. T. Barron, R.
Ramamoorthi, and R. Ng, ``NeRF: Representing scenes as neural radiance
fields for view synthesis,'' in \emph{Proc. ECCV}, 2020, pp.~405-421.

{[}3{]} B. Kerbl, G. Kopanas, T. Leimkühler, and G. Drettakis, ``3D
Gaussian Splatting for real-time radiance field rendering,'' \emph{ACM
Trans. Graph.}, vol.~42, no. 4, pp.~1-14, 2023.

{[}3{]} B. Huang, Y. Yu, D. Chen, et al., ``2D Gaussian Splatting for
geometrically accurate radiance fields,'' in \emph{Proc. SIGGRAPH},
2024.

{[}3{]} M. A. Arshad, T. J. Maxwell, S. A. K. M. Abir, C. W. Bacas, D.
S. E. Moualeu, and A. L. Stutts, ``Evaluating Neural Radiance Fields for
3D plant geometry in field conditions,'' \emph{Plant Phenomics}, vol.~6,
p.~0235, 2024.

{[}3{]} E. Hamuda, M. Glavin, and E. Jones, ``A survey of image
processing techniques for plant extraction and segmentation in the
field,'' \emph{Comput. Electron. Agric.}, vol.~125, pp.~184-199, 2016.

{[}3{]} J. Li, X. Qi, S. H. Nabaei, M. Liu, D. Chen, X. Zhang, X. Yin,
and Z. Li, ``A survey on 3D reconstruction techniques in plant
phenotyping: From classical methods to Neural Radiance Fields (NeRF), 3D
Gaussian Splatting (3DGS), and beyond,'' \emph{Plant Phenomics}, vol.~7,
no. 4, 2025. arXiv:2505.00737

{[}3{]} T. Choi, S. Lee, J. Park, et al., ``NeRF-based 3D reconstruction
pipeline for acquisition and analysis of tomato crop morphology,''
\emph{Front. Plant Sci.}, vol.~15, p.~1439086, 2024.

{[}3{]} Y. Zhang, X. Liu, H. Wang, et al., ``Wheat3DGS: In-field wheat
head reconstruction and phenotyping with 3D Gaussian Splatting,'' in
\emph{Proc. CVPR Workshop on Vision for Agriculture}, 2025.

{[}3{]} J. Shen, L. Chen, Y. Wang, et al., ``PlantGaussian: 3D Gaussian
Splatting for cross-time and cross-scene plant visualization,''
\emph{The Crop Journal}, 2025.

{[}3{]} S. Akhtar, M. F. Shahid, A. Raza, et al., ``Unlocking plant
secrets: A systematic review of 3D imaging in plant phenotyping
techniques,'' \emph{Comput. Electron. Agric.}, vol.~222, p.~109033,
2024.

{[}3{]} A. Guedon and V. Lepetit, ``SuGaR: Surface-Aligned Gaussian
Splatting for efficient 3D mesh reconstruction and high-quality mesh
rendering,'' in \emph{Proc. CVPR}, 2024.

{[}3{]} A. Kirillov, E. Mintun, N. Ravi, et al., ``Segment Anything,''
in \emph{Proc. ICCV}, 2023.

{[}3{]} N. Ravi, V. Gabeur, Y.-T. Hu, et al., ``SAM 2: Segment Anything
in images and videos,'' arXiv:2408.00714, 2024.

{[}3{]} J. W. Abe, J. Ilao, and G. Foliente, ``Promptable leaf
segmentation in plant phenotyping: Research perspectives and
challenges,'' in \emph{Proc. 30th Int. Conf. Mechatronics and Machine
Vision in Practice (M2VIP)}, 2024. doi: 10.1109/M2VIP62436.2024.10745998

{[}3{]} B. Curless and M. Levoy, ``A volumetric method for building
complex models from range images,'' in \emph{Proc. SIGGRAPH}, 1996,
pp.~303-312.

{[}3{]} W. E. Lorensen and H. E. Cline, ``Marching Cubes: A high
resolution 3D surface construction algorithm,'' in \emph{Proc.
SIGGRAPH}, 1987, pp.~163-169.

{[}3{]} S. Xiao, J. Zhang, Y. Liu, et al., ``ICFMNet: Automated
segmentation and 3D phenotypic analysis pipeline for wheat,''
\emph{Comput. Electron. Agric.}, vol.~239, 2025.

{[}3{]} R. Zhang, P. Isola, A. A. Efros, E. Shechtman, and O. Wang,
``The unreasonable effectiveness of deep features as a perceptual
metric,'' in \emph{Proc. CVPR}, 2018, pp.~586-595.

{[}3{]} J. Geng, ``Structured-light 3D surface imaging: a tutorial,''
\emph{Adv. Opt. Photon.}, vol.~3, no. 2, pp.~128-160, 2011.

{[}3{]} M. P. Pound, A. P. French, J. A. Atkinson, D. M. Wells, M. J.
Bennett, and T. P. Pridmore, ``RootNav: Navigating images of complex
root architectures,'' \emph{Plant Physiol.}, vol.~162, no. 4,
pp.~1802-1814, 2013.

{[}3{]} G. Sun and X. Wang, ``Three-dimensional point cloud
reconstruction and morphology measurement method for greenhouse tomato
plants based on RGB-D camera,'' \emph{Comput. Electron. Agric.},
vol.~197, p.~106922, 2022.

{[}3{]} C. Y. Kuo, C. L. Chang, and Y. C. Tsai, ``Multi-view stereo for
plant 3D reconstruction: a comparative study,'' \emph{Biosyst. Eng.},
vol.~216, pp.~198-213, 2022.

{[}3{]} Z. Yang, L. Chen, J. Sun, et al., ``PanicleNeRF: low-cost
high-precision 3D reconstruction and phenotyping of rice panicles with
smartphone,'' \emph{Plant Phenomics}, vol.~6, p.~0279, 2024.

{[}3{]} S. Chopra, R. Khosla, and P. S. Thenkabail, ``AgriNeRF: Neural
Radiance Fields for agricultural scenes under challenging lighting,''
arXiv:2409.15487, 2024.

{[}3{]} Y. Chen, H. Zhang, and W. Li, ``High-fidelity 3D reconstruction
of peach orchards using a 3DGS-Ag model,'' \emph{Comput. Electron.
Agric.}, 2025.

{[}3{]} R. Reena, J. H. Doonan, and Y. H. Liu, ``Wheat3D PartNet:
Annotated 3D point cloud dataset for wheat organ segmentation,''
\emph{Comput. Electron. Agric.}, vol.~238, 2025.

{[}3{]} Z. Gao and X. Su, ``Three-dimensional reconstruction of densely
planted rice seedlings based on multi-view images,'' \emph{Plant
Phenomics}, 2025.

{[}3{]} H. Jiang, X. Sun, S. Li, et al., ``Plant stem occlusion
inpainting with deep reinforcement learning for tomato 3D phenotyping,''
\emph{Comput. Electron. Agric.}, vol.~237, 2025.

{[}3{]} B. Guillard, F. Stella, and P. Fua, ``MeshUDF: Fast and
differentiable meshing of unsigned distance field networks,'' in
\emph{Proc. ECCV}, 2022.

{[}3{]} Y. Furukawa and J. Ponce, ``Accurate, dense, and robust
multiview stereopsis,'' \emph{IEEE Trans. Pattern Anal. Mach. Intell.},
vol.~32, no. 8, pp.~1362-1376, 2010.

{[}3{]} S. Pertuz, D. Puig, and M. A. Garcia, ``Analysis of focus
measure operators for shape-from-focus,'' \emph{Pattern Recognit.},
vol.~46, no. 5, pp.~1415-1432, 2013.

{[}3{]} R. Ferzli and L. J. Karam, ``A no-reference objective image
sharpness metric based on the notion of just noticeable blur (JNB),''
\emph{IEEE Trans. Image Process.}, vol.~18, no. 4, pp.~717-728, 2009.

{[}3{]} J. Park, Y.-W. Tai, D. Cho, and I. S. Kweon, ``A unified
approach of multi-scale deep and hand-crafted features for defocus
estimation,'' in \emph{Proc. CVPR}, 2017, pp.~1736-1745.

{[}3{]} S. Haner and A. Heyden, ``Covariance propagation and next best
view planning for 3D reconstruction,'' in \emph{Proc. ECCV}, 2012.

{[}3{]} C. Mostegel, M. Rumpler, F. Fraundorfer, and H. Bischof,
``UAV-based autonomous image acquisition with multi-view stereo quality
assurance,'' in \emph{Proc. CVPR Workshops}, 2016.

{[}3{]} D. M. Woebbecke, G. E. Meyer, K. Von Bargen, and D. A.
Mortensen, ``Color indices for weed identification under various soil,
residue, and lighting conditions,'' \emph{Trans. ASAE}, vol.~38, no. 1,
pp.~259-269, 1995.

{[}3{]} N. Otsu, ``A threshold selection method from gray-level
histograms,'' \emph{IEEE Trans. Syst. Man Cybern.}, vol.~9, no. 1,
pp.~62-66, 1979.

{[}3{]} A. M. Jimenez, F. Aznarte, and D. A. Riano, ``A review of
computer vision for plant phenotyping in agriculture,'' \emph{Precis.
Agric.}, vol.~24, pp.~1195-1223, 2023.

{[}3{]} M. Minervini, A. Fischbach, H. Scharr, and S. A. Tsaftaris,
``Finely-grained annotated datasets for image-based plant phenotyping,''
\emph{Pattern Recognit. Lett.}, vol.~81, pp.~80-89, 2016.

{[}3{]} M. P. Pound, J. A. Atkinson, A. J. Townsend, et al., ``Deep
machine learning provides state-of-the-art performance in image-based
plant phenotyping,'' \emph{GigaScience}, vol.~6, no. 10, pp.~1-11, 2017.

{[}3{]} D. Cai, Y. Liu, Z. Chen, et al., ``Performance evaluation of
Segment Anything Model for weed detection in cotton fields,''
\emph{Smart Agric. Technol.}, vol.~7, p.~100416, 2024.

{[}3{]} Y. Lu, S. Young, H. Wang, and N. Wijewardane, ``Robust plant
segmentation of proximal aerial images by the Segment Anything Model,''
\emph{Comput. Electron. Agric.}, vol.~218, p.~108715, 2024.

{[}3{]} Z. Fan, K. Wang, K. Wen, Z. Zhu, D. Xu, and Z. Wang,
``LightGaussian: Unbounded 3D Gaussian compression with 15x reduction
and 200+ FPS,'' in \emph{Proc. NeurIPS}, 2024.

{[}3{]} J. C. Lee, D. Rho, X. Sun, J. H. Ko, and E. Park, ``Compact 3D
Gaussian representation for radiance field,'' in \emph{Proc. CVPR},
2024.

{[}3{]} W. Liu, Z. Wu, J. Cao, et al., ``EfficientGS: Streamlining
Gaussian Splatting for large-scale scene reconstruction,''
arXiv:2404.12778, 2024.

{[}3{]} T. Chen, S. Kornblith, M. Norouzi, and G. Hinton, ``A simple
framework for contrastive learning of visual representations,'' in
\emph{Proc. ICML}, 2020, pp.~1597-1607.

{[}3{]} C. G. Northcutt, A. Athalye, and J. Mueller, ``Pervasive label
errors in test sets destabilize machine learning benchmarks,'' in
\emph{Proc. NeurIPS}, 2021.

{[}3{]} J. T. Barron, B. Mildenhall, D. Verbin, P. P. Srinivasan, and P.
Hedman, ``Mip-NeRF 360: Unbounded anti-aliased neural radiance fields,''
in \emph{Proc. CVPR}, 2022, pp.~5470-5479.

\begin{center}\rule{0.5\linewidth}{0.5pt}\end{center}

\emph{Manuscript draft v1.0 --- 2026-05-24. Prepared for Computers and
Electronics in Agriculture.} \emph{Figures and tables use placeholders
with detailed specifications. See accompanying Figure/Table Requirements
Document.}
\end{document}
